\documentclass{article}
\usepackage{../../math_template}

\author{Jonathan Kasongo}
\title{Noether Sheet 7}

\begin{document}
\maketitle
\centerline{
These were my solutions to the Noether Mentoring Scheme Sheet 7. Problem
statements are abridged.
}

\begin{problem}{Problem 1}
Show that for real numbers $x,y$ and $z$ we have
$$
x(x-y) + y(y-z) + z(z-x) \geq 0.
$$
Find all cases where equality occurs.
\end{problem}

\begin{solution}{Solution}
We want to show
$
x^2 + y^2 + z^2 \geq xy + yz + zx.
$
\vspace{0.2cm}

\textit{Theorem --- Cauchy-Schwarz Inequality:
For vectors $\mathbf{u}, \mathbf{v} \in \mathbb{R}^n$, we have
$$
\mathbf{u} \cdot \mathbf{v} \leq \| u \| \| v \|.
$$
with equality if and only if the two vectors are parallel.
\vspace{0.2cm}
}

\textit{Proof:} It follows very simply from the formula for the dot product
$$
\mathbf{u} \cdot \mathbf{v} = \| u \| \| v \| \cos \theta
$$
where $\theta$ is the angle between the two vectors. Now since
$\cos \theta \leq 1$ we have $\mathbf{u} \cdot \mathbf{v} \leq
\| u \| \| v \|$. The equality case happens only when $\cos \theta = 1$
which means $\theta = 0^\circ$, that is the two vectors are parallel.
$\square$
\\

Now let us apply Cauchy-Schwarz with the vectors
$\mathbf{u} = \langle x,y,z \rangle$ and
$\mathbf{v} = \langle y,z,x \rangle$. We have
$$
xy + yz + zx \leq \sqrt{x^2 + y^2 + z^2}\sqrt{y^2 + z^2 + x^2} =
x^2 + y^2 + z^2
$$

which is the desired inequality. The equality case occurs if
$\mathbf{u}$ and $\mathbf{v}$ are parallel which happens only if
\[
\begin{aligned}
\frac{x}{y} = \frac{y}{z} = \frac{z}{x}, &\quad (x,y,z \text{ all non-zero})
\\
(x,y,z) = (0,0,0), &\quad (\text{otherwise}).
\end{aligned}
\]
with the second case being due to the following arguement. Suppose WLOG that
$z=0$. Since the third component of $\mathbf{u}$ and $\mathbf{v}$ have to
be in constant ratio we require $x = 0$. Since the second components
have to be in constant ratio this forces $y=0$, and so all three variables
must be 0. $\blacksquare$
\\

\textit{Remark: This is just a special case of the Re-arrangement
inequality.}
\\
\hline
\vspace{0.2cm}

\textit{Addendums, (after having read solutions):} I believe that my
solution is completely correct, because actually, for $x,y,z$ all non-zero
the following is true
$$
\frac{x}{y} = \frac{y}{z} = \frac{z}{x} \iff x=y=z
$$
like how it was presented in the official solutions.
\end{solution}
\newpage

\begin{problem}{Problem 2}
Three circles touch each other and a line as shown. Prove that
the radii $a,b,c$ where $c$ is the smallest, satisfy the equation
$$
\frac{1}{\sqrt{c}} = \frac{1}{\sqrt{a}} + \frac{1}{\sqrt{b}}
$$
\begin{center}
\begin{asy}
/* Geogebra to Asymptote conversion, documentation at artofproblemsolving.com/Wiki go to User:Azjps/geogebra */
import graph; size(220);
real labelscalefactor = 0.5; /* changes label-to-point distance */
pen dps = linewidth(0.7) + fontsize(10); defaultpen(dps); /* default pen style */
pen dotstyle = black; /* point style */
real xmin = 0, xmax = 18.437582914563066, ymin = 0.42587666093528126, ymax = 11.202016393004262;  /* image dimensions */

pair A = (6.,6.), B = (13.745966692414834,4.), C = (10.36895003862225,1.9526249034443734);
/* draw figures */
draw(circle(A, 5.));
draw(circle(B, 3.));
draw((xmin, 0.*xmin + 1.)--(xmax, 0.*xmax + 1.)); /* line */
draw(circle(C, 0.9526249034443734));
/* dots and labels */
dot(A,dotstyle);
label("$A$", (6.051974368450486,6.1159683804470415), NE * labelscalefactor);
dot(B,dotstyle);
label("$B$", (13.795814691607672,4.122624147497062), NE * labelscalefactor);
dot(C, dotstyle);
label("$C$", (10.413169932662255,2.0447137955734465), NE * labelscalefactor);
clip((xmin,ymin)--(xmin,ymax)--(xmax,ymax)--(xmax,ymin)--cycle);
/* end of picture */
\end{asy}
\end{center}

\end{problem}

\begin{solution}{Solution}
Let the circles have centres $A,B,C$ with radii $a,b,c$ respectively.
The idea is to draw in the following lines,

\begin{center}
\begin{asy}
/* Geogebra to Asymptote conversion, documentation at artofproblemsolving.com/Wiki go to User:Azjps/geogebra */
import graph; size(290);
real labelscalefactor = 0.5; /* changes label-to-point distance */
pen dps = linewidth(0.7) + fontsize(10); defaultpen(dps); /* default pen style */
pen dotstyle = black; /* point style */
real xmin = -0.17988219283588813, xmax = 18.534193727248883, ymin = 0.05699551834098509, ymax = 12.997271270833753;  /* image dimensions */
pen qqwuqq = rgb(0.,0.39215686274509803,0.);
pair A = (6.,6.), B = (13.745966692414834,4.), C = (10.36895003862225,1.9526249034443734), D = (13.745966692414834,1.9526249034443734), F = (6.,1.), G = (13.745966692414834,1.);

filldraw((6.512901162276427,1.)--(6.512901162276427,1.512901162276427)--(6.,1.5129011622764272)--F--cycle, qqwuqq + opacity(0.10000000149011612),  qqwuqq);
filldraw((13.745966692414834,1.5129011622764272)--(13.233065530138408,1.5129011622764272)--(13.233065530138408,1.)--G--cycle, qqwuqq + opacity(0.10000000149011612),  qqwuqq);
/* draw figures */
draw(circle(A, 5.));
draw(circle(B, 3.));
draw((xmin, 0.*xmin + 1.)--(xmax, 0.*xmax + 1.)); /* line */
draw(circle(C, 0.9526249034443734));
draw(C--B);
draw(B--D);
draw(C--D);
draw(C--A);
draw(C--(6.,1.9526249034443734));
draw(A--(6.,1.9526249034443734));
draw((6.,1.9526249034443734)--F,  linetype("4 4"));
draw(D--G,  linetype("4 4"));
/* dots and labels */
dot(A,dotstyle);
label("$A$", (6.058143113859035,6.149950469066279), NE * labelscalefactor);
dot(B,dotstyle);
label("$B$", (13.804900122638314,4.147979556685111), NE * labelscalefactor);
dot(C, dotstyle);
label("$C$", (10.424760828545482,2.0734734663191183), NE * labelscalefactor);
dot(D, dotstyle);
label("$D$", (13.804900122638314,2.0734734663191183), NE * labelscalefactor);
dot((6.,1.9526249034443734), dotstyle);
label("$E$", (6.058143113859035,2.0734734663191183), NE * labelscalefactor);
dot(F, dotstyle);
label("$F$", (6.058143113859035,1.116009116919429), NE * labelscalefactor);
dot(G, dotstyle);
label("$G$", (13.804900122638314,1.116009116919429), NE * labelscalefactor);
clip((xmin,ymin)--(xmin,ymax)--(xmax,ymax)--(xmax,ymin)--cycle);
/* end of picture */
\end{asy}
\end{center}

First look at $\triangle CBD$. Side $BC = b+c$ since the two circles
are tangent to each other. Side
$BD = BG - BD = b - c$. Thus by Pythagoras
$CD = \sqrt{(b+c)^2 - (b-c)^2} = 2\sqrt{bc}$. \vspace{0.2cm}

Similarly look at $\triangle CEA$. Side $AC = a+c$ since the two circles
are tangent. Side $AE = AF - EF = a - c$. Thus by Pythagoras
$CE = \sqrt{(a+c)^2 - (a-c)^2} = 2\sqrt{ac}$.
So putting these two together, $FG = ED = 2\sqrt{bc} + 2\sqrt{ac}$.
Now consider the following

\begin{center}
\begin{asy}
/* Geogebra to Asymptote conversion, documentation at artofproblemsolving.com/Wiki go to User:Azjps/geogebra */
import graph; size(290);
real labelscalefactor = 0.5; /* changes label-to-point distance */
pen dps = linewidth(0.7) + fontsize(10); defaultpen(dps); /* default pen style */
pen dotstyle = black; /* point style */
real xmin = -0.17988219283588813, xmax = 18.337007966034857, ymin = 0.55699551834098509, ymax = 12.097271270833753;  /* image dimensions */
pen qqwuqq = rgb(0.,0.39215686274509803,0.);
pair A = (6.,6.), B = (13.745966692414834,4.), C = (10.36895003862225,1.9526249034443734), D = (13.745966692414834,1.9526249034443734), F = (6.,1.), G = (13.745966692414834,1.), H = (6.,4.);

filldraw((6.512901162276427,1.)--(6.512901162276427,1.5129011622764272)--(6.,1.5129011622764272)--F--cycle, qqwuqq + opacity(0.10000000149011612),  qqwuqq);
filldraw((13.745966692414834,1.5129011622764272)--(13.233065530138408,1.5129011622764272)--(13.233065530138408,1.)--G--cycle, qqwuqq + opacity(0.10000000149011612),  qqwuqq);
/* draw figures */
draw(circle(A, 5.));
draw(circle(B, 3.));
draw((xmin, 0.*xmin + 1.)--(xmax, 0.*xmax + 1.)); /* line */
draw(circle(C, 0.9526249034443734));
draw(B--D);
draw(A--(6.,1.9526249034443734));
draw((6.,1.9526249034443734)--F,  linetype("4 4"));
draw(D--G,  linetype("4 4"));
draw(A--B);
draw(H--B);
/* dots and labels */
dot(A,dotstyle);
label("$A$", (6.058143113859036,6.149950469066279), NE * labelscalefactor);
dot(B,dotstyle);
label("$B$", (13.804900122638315,4.147979556685111), NE * labelscalefactor);
dot(C, dotstyle);
label("$C$", (10.192648258994044,2.0734734663191183), NE * labelscalefactor);
dot(D, dotstyle);
label("$D$", (13.804900122638315,2.0734734663191183), NE * labelscalefactor);
dot((6.,1.9526249034443734), dotstyle);
label("$E$", (6.058143113859036,2.0734734663191183), NE * labelscalefactor);
dot(F, dotstyle);
label("$F$", (6.058143113859036,1.116009116919429), NE * labelscalefactor);
dot(G, dotstyle);
label("$G$", (13.804900122638315,1.116009116919429), NE * labelscalefactor);
dot(H, dotstyle);
label("$H$", (6.058143113859036,4.118965485491181), NE * labelscalefactor);
clip((xmin,ymin)--(xmin,ymax)--(xmax,ymax)--(xmax,ymin)--cycle);
/* end of picture */
\end{asy}
\end{center}

Look at $\triangle AHB$. Side $AB = a+b$ since the two circles are tangent.
Side $AH = AF - HF = a - b$. So by Pythagoras,
$HB = \sqrt{(a+b)^2 - (a-b)^2} = 2\sqrt{ab}$. That is also equal to $FG$.
So we deduce
$$
2\sqrt{ab} = 2\sqrt{bc} + 2\sqrt{ac}
$$
and dividing by $2\sqrt{abc}$ yields the desired result,
$$
\frac{1}{\sqrt{c}} = \frac{1}{\sqrt{a}} + \frac{1}{\sqrt{b}}. \ \blacksquare
$$
\\

\hline
\vspace{0.2cm}

\textit{Addendums (after having read solutions):} A classical result that
comes up fairly often usually in multiple choice type math contests.
I have seen this idea before.
\end{solution}
\newpage

\begin{problem}{Problem 3}
Find all quadruples $(a,b,c,d)$ of positive integers such that
$a \leq b \leq c \leq d$ and
$$
a + b + c + d = ab + cd.
$$
\end{problem}

\begin{solution}{Solution}
Due to this ordering condition we notice that
$$
a+b+c+d \leq d+d+d+d = 4d.
$$
But since this sum is the same as $ab+cd$ we note that
$ab+cd \leq 4d$. Now if $c \geq 4$ the quantity $ab+cd$ will be too large,
so $c < 3$. \vspace{0.2cm}

If $c = 3$, we have
\[
\begin{aligned}
a + b + 3 + d &= 3d + ab \\
a + b - ab + 3 &= 2d \\
\end{aligned}
\]
If $b = 3$ then
\[
\begin{aligned}
-2a + 6 &= 2d \\
-a + 3 &= d \\
\end{aligned}
\]
but this implies that $d < 3 = c$ which is a contradiction.
\vspace{0.2cm}

If $b = 2$ then
\[
\begin{aligned}
-a + 5 &= 2d \\
\frac{5-a}{2} &= d
\end{aligned}
\]
but this implies that $d < 2.5 < c$, a contradiction.
\vspace{0.2cm}

If $b = 1$, then this forces $a=1$ due to the ordering condition.
We have
\[
\begin{aligned}
1 + 2 + 3 + d &= 2 + 3d \\
4 &= 2d
\end{aligned}
\]
so $d=2$. But $d < 3 = c$ is a contradiction.
\vspace{0.2cm}

If $c = 2$, we have
\[
\begin{aligned}
a + b + 2 + d &= ab + 2d \\
a + b - ab + 2 &= d
\end{aligned}
\]
If $b = 2$ then
$$
-a + 4 = d
$$
If $a = 1$ then $d = 3$ which yields a solution
$(a,b,c,d) = (1,2,2,3)$. \vspace{0.2cm}

If $a = 2$ then $d = 2$ which yields a solution
$(a,b,c,d) = (2,2,2,2)$. \vspace{0.2cm}

If $a > 2$ then $d < 2 = c$ which is a contradiction.
\vspace{0.2cm}

If $c = 1$, we have $a=b=c=1$ due to the ordering condition. This gives
$3 + d = 1 + d$ which clearly has no solutions. \vspace{0.2cm}

Thus the only solutions are $(a,b,c,d) = (1,2,2,3), (2,2,2,2)$.
$\blacksquare$
\\

\hline
\vspace{0.2cm}

\textit{Addendums (after having read solutions):} On this problem I simply
just forgot the case if $c=2$ and $b=1$ and this yields the final solution
$(1,1,2,3)$.
\end{solution}
\newpage

\begin{problem}{Problem 4}
(a) How many different values is it possible to make using three different
British coins? \vspace{0.2cm}

(b) In how many ways can three numbers be chosen from the set
$\{1, 2, 3, 4, 5, 6, 7, 8, 9, 10\}$ such that no two of them are
consecutive? \vspace{0.2cm}

(c) You should notice that these two numbers are the same. Why is this?
\end{problem}

\begin{hint}{Preliminaries}
Preliminaries: Let $f : A \rightarrow B$ be a
function. The set of inputs $A$ is the domain of $f$. The set of all
outputs $B$ is the co-domain of $f$ or also known as the image of $f$. I
may use both terminology. \\

An injective function
$f : A \rightarrow B$ is a function with the following property.
$$
f(x_1) = f(x_2) \iff x_1 = x_2
$$
Roughly speaking that means that exactly one unique input can
map to a one unique output. For example $f(x) = \ln(x)$ is injective
whilst $f(x) = x^2$ isn't. \\

A surjective function is a function
$f : A \rightarrow B$ is a function with the following property.
$$
\text{For every } b \in B, \text{ there exists some } a \in A
\text{ such that } f(a) = b
$$
Roughly speaking that means that the function hits every element in the
co-domain. For example $f(x) = 2x$ is surjective whilst $f(x) = e^x$
isn't (if we assume the codomain is the entirety of $\mathbb{R}$). \\

A bijective function is a function $f : A \rightarrow B$ which is both
injective and surjective. Bijective functions are also known as "one to one"
functions because every unique element in the domain $A$ is paired up
with it's own unique element in the image $B$. I will skip out the
formal definition here (because im lazy) but the following diagram gives
us a good intuition for the idea of a bijective function.

\begin{center}
\begin{asy}
import graph;
size(150, 150);

// Create the containers for the sets
draw(ellipse((0,0), 0.8, 1.5), linewidth(1pt));
draw(ellipse((4,0), 0.8, 1.5), linewidth(1pt));

// Set Labels
label("Set $A$", (0, 1.8), N);
label("Set $B$", (4, 1.8), N);

// Define coordinates for the points
pair[] setA = {(0, 0.8), (0, 0), (0, -0.8)};
pair[] setB = {(4, 0.8), (4, 0), (4, -0.8)};

// Labels for elements
string[] labelsA = {"$a_1$", "$a_2$", "$a_3$"};
string[] labelsB = {"$b_1$", "$b_2$", "$b_3$"};

// Draw dots and labels
for(int i=0; i<3; ++i) {
    dot(setA[i]);
    label(labelsA[i], setA[i], W);
    dot(setB[i]);
    label(labelsB[i], setB[i], E);
}

// Draw the Bijective Mapping (One-to-One and Onto)
// We'll scramble them slightly to show it doesn't have to be parallel lines
draw(setA[0] -- setB[1], Arrow(HookHead));
draw(setA[1] -- setB[0], Arrow(HookHead));
draw(setA[2] -- setB[2], Arrow(HookHead));

// Label the function
label("$f$", (2, 0.6), N);
\end{asy}
\end{center}

\end{hint}
\newpage

\begin{solution}{Solution}
(a) The answer is just $\binom{8}{3} = 56$ since each triple of coins
produces it's own unique sum. We prove this statement by showing that
$f(x,y,z)$ is an injection, where $f$ is defined below. \vspace{0.2cm}

The eight British coins have values: 1p, 2p, 5p, 10p, 20p, 50p, £1, £2.
It is nicer to think about these values all in terms of pence. Define the
sequence
$$
(a_n)_{n=1}^8 = (1, 2, 5, 10, 20, 50, 100, 200).
$$

For any triple of integers $(x,y,z)$ with $1 \leq x < y < z \leq 8$,
define $f(x,y,z) = a_x + a_y + a_z$.
\\

\textit{Claim 1: Let $2 \leq n \leq 8$ be an integer. Then
$
a_1 + a_2 + \cdots + a_{n-1} < a_n.
$
}\vspace{0.2cm}

\textit{Proof: } Check each case manually. $\square$ \\

\textit{Claim 2: The function $f$ is injective.} \vspace{0.2cm}

\textit{Proof: } Suppose FTSOC that $f$ is not injective. Then there exists
some triple $(x,y,z)$ with the following property. There exists some
different triple $(u,v,w)$ such that $f(x,y,z) = f(u,v,w)$. Now aiming for
a contradiction we are going to show that $(x,y,z)$ and $(u,v,w)$ are in
fact the same triple. \vspace{0.2cm}

First we show that $w = z$, by showing that $w > z$ and $w < z$ both give
contradictions. Note that $3\leq w,z \leq 8$ due to ordering conditions.
\begin{itemize}
\item Suppose that $w > z$. Because claim 1 tells us that any sum $S$ of distinct
elements that come before $a_w$ in the sequence, satisfies $S < a_w$, we
have $f(x,y,z) < a_w < f(u,v,w)$ a contradiction.

\item
Suppose that $w < z$. Similarly because claim 1 tells us that any sum $S$
of distinct elements that come before $a_z$ in the sequence,
satisfies $S < a_z$, we have $f(u,v,w) < a_z < f(x,y,z)$ a contradiction.

\end{itemize}

Thus the only possibility is that $w = z$. Now we are going to show that
$v = y$ using the same strategy. Because $w=z$ our initial assumption
reduces to $a_u + a_v = a_x + a_y$. Note that $2 \leq v,y \leq 7$ due to
ordering conditions.

\begin{itemize}
\item Suppose that $v > y$. Because claim 1 tells us that any sum $S$ of
distinct elements that come before $a_v$ in the sequence satisfies
$S < a_v$, we have $a_x + a_y < a_v < a_u + a_v$ a contradiction.
\item Suppose that $v < y$. Because claim 1 tells us that any sum $S$ of
distinct elements that come before $a_y$ in the sequence, satisfies
$S < a_y$, we have $a_u + a_v < a_y < a_x + a_y$ a contradiction.
\end{itemize}

Thus the only possibility is that $v = y$. So now our initial assumption
reduces to $a_u = a_x$. Since there are no repeated values in $(a_n)$ it
must be the case that $u = v$ and so in fact $(x,y,z)$ and $(u,v,w)$
are the same exact triple, contradiction. $\square$
\\

Now due to claim 2, each selection of three distinct British coins
produces a unique sum, so the number of different values one can make is
just the number of unordered triples of three coins without repitition and
that works out to be
$$
\binom{8}{3} = 56. \ \square
$$
\\

(b) We will use complementary counting. Call an unordered triple $(c,b,a)$
of integers $1 \leq c < b < a \leq 10$ \textit{nice}.
\vspace{0.2cm}

We want to count
the number $N$ of nice triples that have at least one pair of consecutive
numbers. For each pair $(n, n+1)$, $(1 \leq n \leq 9)$ there are 8 other
numbers that we can put in the triple giving $9(8)$ such nice triples.
However we have double counted all the nice triples of the form
$(n, n+1, n+2)$, $(1 \leq n \leq 8)$ since they have two pairs of
consecutive numbers. So in fact $N = 9(8) - 8 = 64$.
\vspace{0.2cm}

Thus the desired count works out as
$$
\binom{10}{3} - 64 = 56. \ \square
$$
\\

(c) Call a nice triple \textit{super nice} if the triple has
no pair of consecutive numbers. Call a nice triple $(x,y,z)$
\textit{almost nice} if it satisfies $(2\leq x < y < z \leq 9)$. There
are exactly $\binom{8}{3}$ almost nice triples. \vspace{0.2cm}

I am going to show that
there is a bijection (a "perfect" one to one correspondence)
between super nice
triples and almost nice triples, and this will force our conclusion
"The number of super nice triples is the same as the number of almost nice
triples. This is why there are $\binom{8}{3}$ super nice triples."
\vspace{0.2cm}

Let $f$ be the bijection. Precisely we have the following
$$
f : \{ \text{Super nice triples} \} \rightarrow
\{ \text{Almost nice triples} \}.
$$
Now we will describe the map $f$. Consider some arbitrary super nice
triple $(p,q,r)$. Note that $3\leq q\leq 8$ due to the ordering
restrictions on super nice triples.
We will map it to the almost nice triple $(x,q,z)$ with
the following rules.
\begin{itemize}
\item If $p = 1$ then $x = q - 1$. Otherwise $x = p$.
\item If $q = 10$ then $z = q + 1$. Otherwise $z = r$.
\end{itemize}
With the inequality $3\leq q \leq 8$ it can be checked that the output
triple $(x,q,z)$ is indeed almost nice. Since we have found a bijection
between super nice triples and almost nice triples, the number of triples
of both forms is the exact same and this is why there are
$\binom{8}{3} = 56$ super nice triples. $\blacksquare$
\\

\hline
\vspace{0.2cm}

\textit{Addendums (after having read solutions):} I believe that
my solutions to (a) and (b) are correct and justified. The hard part is
part (c), it's quite difficult I feel to communicate one's solution clearly
for this one. \vspace{0.2cm}

I hope that you can at least understand the gist of my solution, if I'm
being perfectly honest I didn't really understand the official solution here
I feel that it lacks a lot of rigour.}
\end{solution}
\newpage

\begin{problem}{Problem 5}
An acute-angled triangle $ABC$ is inscribed inside a circle. The altitudes
of the triangle $AD, BE, CF$ meet at $H$, the orthocentre of the triangle.
\vspace{0.2cm}

(a) If line $AD$ is extended to meet the circle at $P$, prove that
$HD = DP$.
\vspace{0.2cm}

(b) Prove that $H$ is the incentre of triangle $DEF$.
\end{problem}

\begin{solution}{Solution}
(a) Here's a diagram.

\begin{center}
\begin{asy}
/* Geogebra to Asymptote conversion, documentation at artofproblemsolving.com/Wiki go to User:Azjps/geogebra */
import graph; size(230);
real labelscalefactor = 0.5; /* changes label-to-point distance */
pen dps = linewidth(0.7) + fontsize(10); defaultpen(dps); /* default pen style */
pen dotstyle = black; /* point style */
real xmin = -0.3, xmax = 15.2, ymin = -3.2, ymax = 12.207635095338665;  /* image dimensions */
pen ttttff = rgb(0.2,0.2,1.); pen qqzzff = rgb(0.,0.6,1.);
pair A = (1.,1.), B = (14.,1.), C = (3.,10.), D = (6.212871287128713,7.371287128712871), F = (3.,1.), H = (3.,3.4444444444444446), P = (9.425742574257425,11.298129812981298);

filldraw(circle((7.5,4.277777777777778), 7.279685924577642), ttttff + opacity(0.05000000074505806),  ttttff);
filldraw((5.584895176786731,7.885085764447219)--(5.071096541052382,7.257109654105238)--(5.699072651394364,6.74331101837089)--D--cycle, qqzzff + opacity(0.25),  qqzzff);
filldraw((3.8113834070541017,1.)--(3.8113834070541017,1.8113834070541015)--(3.,1.8113834070541017)--F--cycle, qqzzff + opacity(0.25),  qqzzff);
filldraw((2.403826688696536,3.5769274025118807)--(2.5798404626552434,4.3689893853260635)--(1.7877784798410605,4.5450031592847715)--(1.611764705882353,3.7529411764705882)--cycle, qqzzff + opacity(0.25),  qqzzff);
filldraw(arc(C,1.1474694185404002,-102.52880770915152,-39.289406862500364)--(3.,10.)--cycle, black + opacity(0.0));
filldraw(arc(P,1.1474694185404002,-129.28940686250036,-66.0500060158492)--(9.425742574257425,11.298129812981298)--cycle, black + opacity(0.0));
filldraw(arc(A,1.1474694185404002,410.71059313749964,437.4711922908485)--(1.,1.)--cycle, black + opacity(0.0));
filldraw(arc(B,1.1474694185404002,473.94999398415075,500.71059313749964)--(14.,1.)--cycle, black + opacity(0.0));
filldraw(arc(H,1.1474694185404002,167.47119229084848,230.71059313749964)--(3.,3.4444444444444446)--cycle, black + opacity(0.0));
filldraw(arc(H,1.1474694185404002,347.4711922908485,410.71059313749964)--(3.,3.4444444444444446)--cycle, black + opacity(0.0));
filldraw(arc(B,1.3769633022484804,500.71059313749964,527.4711922908485)--(14.,1.)--cycle, black + opacity(0.0));
/* draw figures */
draw(C--B);
draw(B--A);
draw(C--A);
draw(C--F);
draw(A--D);
draw(B--(1.611764705882353,3.7529411764705882));
draw(D--P);
draw(P--B,  linetype("4 4"));
draw(arc(C,1.1474694185404002,-102.52880770915152,-39.289406862500364));
draw(arc(C,0.9982983941301482,-102.52880770915152,-39.289406862500364));
draw(arc(P,1.1474694185404002,-129.28940686250036,-66.0500060158492));
draw(arc(P,0.9982983941301482,-129.28940686250036,-66.0500060158492));
draw(arc(A,1.1474694185404002,410.71059313749964,437.4711922908485));
draw(arc(A,0.9982983941301482,410.71059313749964,437.4711922908485));
draw(arc(A,0.8491273697198962,410.71059313749964,437.4711922908485));
draw(arc(B,1.1474694185404002,473.94999398415075,500.71059313749964));
draw(arc(B,0.9982983941301482,473.94999398415075,500.71059313749964));
draw(arc(B,0.8491273697198962,473.94999398415075,500.71059313749964));
draw(arc(H,1.1474694185404002,167.47119229084848,230.71059313749964));
draw(arc(H,0.9982983941301482,167.47119229084848,230.71059313749964));
draw(arc(H,1.1474694185404002,347.4711922908485,410.71059313749964));
draw(arc(H,0.9982983941301482,347.4711922908485,410.71059313749964));
draw(arc(B,1.3769633022484804,500.71059313749964,527.4711922908485));
draw(arc(B,1.2277922778382284,500.71059313749964,527.4711922908485));
draw(arc(B,1.0786212534279764,500.71059313749964,527.4711922908485));
/* dots and labels */
dot(A,dotstyle);
label("$A$", (1.1024440394066168,1.2352928484648533), NW * labelscalefactor);
dot(B,dotstyle);
label("$B$", (14.091797857283948,1.2352928484648533), NE * labelscalefactor);
dot(C,dotstyle);
label("$C$", (3.0990408276669132,10.23145308982147), NE * labelscalefactor);
dot(D, dotstyle);
label("$D$", (6.3119551995800345,7.546374650436969), NE * labelscalefactor);
dot((1.611764705882353,3.7529411764705882), dotstyle);
label("$E$", (1.699128137047625,3.943320676220161), NW * labelscalefactor);
dot(F, dotstyle);
label("$F$", (3.0990408276669132,1.189394071723238), NE * labelscalefactor);
dot(H, dotstyle);
label("$H$", (3.0990408276669132,3.6220292390288535), NW * labelscalefactor);
dot(P, dotstyle);
label("$P$", (9.524869571493156,11.470720061845086), NE * labelscalefactor);
clip((xmin,ymin)--(xmin,ymax)--(xmax,ymax)--(xmax,ymin)--cycle);
/* end of picture */
\end{asy}
\end{center}

First let's study $\triangle CDA$.
Call $\angle DCA = \theta$. Then $\angle CAD = 90 - \theta$.
Now let's look at $\triangle EHA$. We have $\angle EHA = \theta$ since
$\angle HEA = 90$. Now due to vertically opposite angles $\angle DHB =
\angle EHA = \theta$.
\vspace{0.2cm}

However due to the "Angles subtending the same segment are equal"
theorem we also have $\angle APB = \angle ACB = \theta$ or equivalently
$\angle ACD = \theta = \angle DPB$. Now notice that
$\angle HDB = \angle PDB = 90$ due to vertically opposite angles +
angles on a straight line sum to 180. Moreover observe that
$\triangle DPB$ and $\triangle DBH$ share two pairs of equal angles, and
since angles in a triangle sum to 180, their third angles are equal also.
Thus the two triangles are similar by AAA criterion. \vspace{0.2cm}

But since the two triangles share a side $DB$, they are actually
congruent to each other by ASA criterion. This triangle congruence
implies the desired result $DH = DP$. $\square$
\\ \bigskip

(b) Here's a diagram. I will show that $CF$ bisects $\angle EFD$, and
then one can use a similar process to show the other bisections.

\begin{center}
\begin{asy}
/* Geogebra to Asymptote conversion, documentation at artofproblemsolving.com/Wiki go to User:Azjps/geogebra */
import graph; size(250);
real labelscalefactor = 0.5; /* changes label-to-point distance */
pen dps = linewidth(0.7) + fontsize(10); defaultpen(dps); /* default pen style */
pen dotstyle = black; /* point style */
real xmin = -2.8, xmax = 12., ymin = 0., ymax = 9.587;  /* image dimensions */
pen ttttff = rgb(0.2,0.2,1.); pen yqqqqq = rgb(0.5019607843137255,0.,0.); pen ffzzqq = rgb(1.,0.6,0.); pen qqzzff = rgb(0.,0.6,1.);
pair A = (1.,1.), B = (11.,1.), C = (4.,9.), D = (6.663716814159292,5.95575221238938), F = (4.,1.), H = (4.,3.625);

filldraw((4.489927763445923,1.)--(4.489927763445923,1.4899277634459227)--(4.,1.4899277634459227)--F--cycle, ttttff + opacity(0.25),  ttttff);
filldraw((2.6916103721443294,4.115646110445876)--(2.863635494575165,4.574379770261438)--(2.404901834759603,4.746404892692274)--(2.232876712328767,4.287671232876712)--cycle, ttttff + opacity(0.25),  ttttff);
filldraw((6.2950085973000265,5.633132522637523)--(6.617628287051883,5.264424305778257)--(6.986336503911149,5.587043995530115)--D--cycle, ttttff + opacity(0.25),  ttttff);
filldraw((2.232876712328767,4.287671232876712)--D--F--cycle, yqqqqq + opacity(0.15000000596046448),  yqqqqq);
filldraw(circle((2.5,5.), 4.272001872658766), ffzzqq + opacity(0.05000000074505806),  ffzzqq);
filldraw(circle((2.5,2.3125), 1.9931523398877469), qqzzff + opacity(0.15000000596046448),  qqzzff);
/* draw figures */
draw(A--B);
draw(B--C);
draw(C--A);
draw(C--F);
draw(A--D);
draw((2.232876712328767,4.287671232876712)--B);
draw((2.232876712328767,4.287671232876712)--D);
draw(D--F);
draw(F--(2.232876712328767,4.287671232876712));
draw((2.232876712328767,4.287671232876712)--D,  yqqqqq);
draw(D--F,  yqqqqq);
draw(F--(2.232876712328767,4.287671232876712),  yqqqqq);
/* dots and labels */
dot(A,dotstyle);
label("$A$", (1.0574064859562335,1.1323965470536064), NE * labelscalefactor);
dot(B,dotstyle);
label("$B$", B, NE * labelscalefactor);
//label("$B$", (9.3856135874893,-0.6551886710791148), NE * labelscalefactor);
dot(C,dotstyle);
label("$C$", (4.050572432597069,9.141886904268436), NE * labelscalefactor);
dot(D, dotstyle);
label("$D$", (6.725021634919668,6.065577459109798), NE * labelscalefactor);
dot((2.232876712328767,4.287671232876712), dotstyle);
label("$E$", (2.2907017139702814,4.402707488753778), NE * labelscalefactor);
dot(F, dotstyle);
label("$F$", (4.050572432597069,1.1046820475476729), NE * labelscalefactor);
dot(H, dotstyle);
label("$H$", (4.050572432597069,3.7375595006113707), NE * labelscalefactor);
clip((xmin,ymin)--(xmin,ymax)--(xmax,ymax)--(xmax,ymin)--cycle);
/* end of picture */
\end{asy}
\end{center}

\textit{Claim 1: $AFDC$ is cyclic.}
\vspace{0.2cm}

\textit{Proof:} Look at $\triangle BDA$. We have
$\angle BDA = 180 - (90 + \angle B) = 90 - \angle B$. Look at
$\triangle FCB$. We have
$\angle FCB = 180 - (90 + \angle B) = 90 - \angle B$. Now since
$\angle BDA = \angle FCB$ both subtend the segement $DF$, by the converse
of the "Angles subtending the same segment are equal" theorem we know that
$A,F,D,C$ are concyclic. $\square$ \\

\textit{Claim 2: $AFHE$ is cyclic.} \vspace{0.2cm}

\textit{Proof:} We have opposite angles summing to 180 since
$\angle AEH = \angle AFH = 90$. These two summing to 180 forces the other
pair of opposite angles to sum to 180 also, because angles in a
quadrilateral sum to 360. Thus $AFHE$ is cyclic. $\square$ \\

By angles subtending the same segment are equal on $AFDC$ we have
$\angle CFD = \angle CAD$. By the same theorem on $AFHE$ we have
$\angle EAH = \angle EFH$, or equivalently $\angle CAD = \angle EFH$.
Now since $\angle CFD = \angle EFH$ the line $CF$
has in fact bisected $\angle EFD$. $\blacksquare$
\\

\hline
\vspace{0.2cm}

\textit{Addendum (after having read solutions):} I think part (a) is fine
here. For part (b) I only showed that $CF$ bisects $\angle EFD$ out
of laziness, I kinda just hoped that the same method will work to show
the two other bisections (I don't see any reason why it should work out
that way).
\end{solution}
\newpage

\begin{problem}{Problem 6}
(a) Suppose twelve $2 \times 1$ dominoes are placed on a $5 \times 5$
chessboard such that each covers exactly two squares. Which squares can
possibly be left uncovered? \vspace{0.2cm}

(b) Suppose instead that eight $3 \times 1$ triominoes are used, each
covering exactly three squares. Which squares can be left uncovered now?
\vspace{0.2cm}

(c) What if six $4 \times 1$ tetrominoes are used, with each covering
four squares. Which squares can be left uncovered in that case?
\end{problem}

\begin{solution}{Solution}
(a) There are two ways that a $5 \times 5$ chessboard could be coloured,
with 12 black squares and 13 white squares or with 13 black squares
and 12 white squares. WLOG let us assume the latter case,

\begin{center}
\begin{asy}
/* Geogebra to Asymptote conversion, documentation at artofproblemsolving.com/Wiki go to User:Azjps/geogebra */
import graph; size(200);
real labelscalefactor = 0.5; /* changes label-to-point distance */
pen dps = linewidth(0.7) + fontsize(10); defaultpen(dps); /* default pen style */
pen dotstyle = black; /* point style */
real xmin = -1.5154308812998338, xmax = 10.888679016798562, ymin = 0.23324824404005604, ymax = 6.780269563036372;  /* image dimensions */


filldraw((1.,6.)--(2.,6.)--(2.,5.)--(1.,5.)--cycle, black + opacity(1.0));
filldraw((3.,6.)--(4.,6.)--(4.,5.)--(3.,5.)--cycle, black + opacity(1.0));
filldraw((5.,6.)--(6.,6.)--(6.,5.)--(5.,5.)--cycle, black + opacity(1.0));
filldraw((2.,5.)--(3.,5.)--(3.,4.)--(2.,4.)--cycle, black + opacity(1.0));
filldraw((4.,5.)--(5.,5.)--(5.,4.)--(4.,4.)--cycle, black + opacity(1.0));
filldraw((1.,4.)--(2.,4.)--(2.,3.)--(1.,3.)--cycle, black + opacity(1.0));
filldraw((3.,4.)--(4.,4.)--(4.,3.)--(3.,3.)--cycle, black + opacity(1.0));
filldraw((5.,4.)--(6.,4.)--(6.,3.)--(5.,3.)--cycle, black + opacity(1.0));
filldraw((2.,3.)--(3.,3.)--(3.,2.)--(2.,2.)--cycle, black + opacity(1.0));
filldraw((4.,3.)--(5.,3.)--(5.,2.)--(4.,2.)--cycle, black + opacity(1.0));
filldraw((5.,2.)--(6.,2.)--(6.,1.)--(5.,1.)--cycle, black + opacity(1.0));
filldraw((3.,2.)--(4.,2.)--(4.,1.)--(3.,1.)--cycle, black + opacity(1.0));
filldraw((1.,2.)--(2.,2.)--(2.,1.)--(1.,1.)--cycle, black + opacity(1.0));
/* draw figures */
draw((1.,6.)--(2.,6.));
draw((2.,6.)--(2.,5.));
draw((2.,5.)--(1.,5.));
draw((1.,5.)--(1.,6.));
draw((3.,6.)--(4.,6.));
draw((4.,6.)--(4.,5.));
draw((4.,5.)--(3.,5.));
draw((3.,5.)--(3.,6.));
draw((5.,6.)--(6.,6.));
draw((6.,6.)--(6.,5.));
draw((6.,5.)--(5.,5.));
draw((5.,5.)--(5.,6.));
draw((2.,5.)--(3.,5.));
draw((3.,5.)--(3.,4.));
draw((3.,4.)--(2.,4.));
draw((2.,4.)--(2.,5.));
draw((4.,5.)--(5.,5.));
draw((5.,5.)--(5.,4.));
draw((5.,4.)--(4.,4.));
draw((4.,4.)--(4.,5.));
draw((1.,4.)--(2.,4.));
draw((2.,4.)--(2.,3.));
draw((2.,3.)--(1.,3.));
draw((1.,3.)--(1.,4.));
draw((3.,4.)--(4.,4.));
draw((4.,4.)--(4.,3.));
draw((4.,3.)--(3.,3.));
draw((3.,3.)--(3.,4.));
draw((5.,4.)--(6.,4.));
draw((6.,4.)--(6.,3.));
draw((6.,3.)--(5.,3.));
draw((5.,3.)--(5.,4.));
draw((2.,3.)--(3.,3.));
draw((3.,3.)--(3.,2.));
draw((3.,2.)--(2.,2.));
draw((2.,2.)--(2.,3.));
draw((4.,3.)--(5.,3.));
draw((5.,3.)--(5.,2.));
draw((5.,2.)--(4.,2.));
draw((4.,2.)--(4.,3.));
draw((5.,2.)--(6.,2.));
draw((6.,2.)--(6.,1.));
draw((6.,1.)--(5.,1.));
draw((5.,1.)--(5.,2.));
draw((3.,2.)--(4.,2.));
draw((4.,2.)--(4.,1.));
draw((4.,1.)--(3.,1.));
draw((3.,1.)--(3.,2.));
draw((1.,2.)--(2.,2.));
draw((2.,2.)--(2.,1.));
draw((2.,1.)--(1.,1.));
draw((1.,1.)--(1.,2.));
draw((1.,6.)--(6.,6.));
draw((6.,6.)--(6.,1.));
draw((6.,1.)--(1.,1.));
draw((1.,1.)--(1.,6.));
/* dots and labels */
clip((xmin,ymin)--(xmin,ymax)--(xmax,ymax)--(xmax,ymin)--cycle);
/* end of picture */
\end{asy}
\end{center}

Each of the twelve $2 \times 1$ dominoes covers exactly one white square
and one black square. All together they will cover 12 white squares and
12 black squares. Thus the one square left uncovered must be black.
Now we are going to show that any black square could be left uncovered.
We show that there is an algorithm that allows us to always cover the board.
\\

\textit{Algorithm 1.}
\vspace{0.2cm}

We are going to tile a $5\times 5$ chessboard with one black square
removed. We are going to dissect the shape into twelve $2 \times 1$ pieces.

\begin{enumerate}
\item Consider the topmost row of the board, now take the leftmost
square on this row, call the square $s$.
\item If there is a square just below $s$ then, remove both $s$ and this
square from the board. Otherwise take the square to the right of $s$ and
remove both $s$ and this square from the board.
\item If the board still has squares left go back to step 1, otherwise
terminate the algorithm here.
\end{enumerate}
\vspace{0.2cm}

Now we are going to prove algorithm 1 always works to tile the board.
\\

\textit{Proof of correctness:} We need to show that the square below or
to the right of $s$ always exists in step 2.
Consider the square $s$ that we refer to in step 2.
The following is a key lemma in our arguement, call it $(*)$.
Observe that since we always take $s$ to be the topmost then the leftmost
square available, the previously discarded squares were either above $s$
or on the left of $s$.
\vspace{0.2cm}

\textit{Case 1: $s$ is black.}
If $s$ is black, and if $s$ doesn't lie on the bottom row,
the square below it $b$, is white. The square $b$ exists because it wasn't
previously discarded due to $(*)$, and
it isn't the one square that was removed from the chessboard, since we
removed a black square from the board. \vspace{0.2cm}

Now if $s$ does lie on the bottom row
then $s$ has a square to the right of it unless $s$ happens to be the
square in the very bottom right of the chessboard. But since $s$ is the
topmost then the leftmost square available, this must
mean that $s$ is the only square left on the entire chessboard at the
moment. However since there were initially
$24 \equiv 0 \ (\text{mod } 2)$ squares, and each time we
are removing $2 \equiv 0 \ (\text{mod } 2)$ squares from the shape, the
number of squares currently part of the shape $N$, satisfies the invariant
$N \equiv 0 \ (\text{mod } 2)$. Thus $s$ cannot be the only square left,
a contradiction.
\vspace{0.2cm}

\textit{Case 2: $s$ is white.}
If $s$ is white, and not on the bottom row, then there is a black square
below it $b$, only if $b$ was not the square we chose to remove from the
board. Suppose that $b$ was in fact the square we chose to remove. If
$s$ doesn't lie on the right most column then we can just pick the black
square to the right of $s$. Suppose that $s$ actually lies on the right most
column. Due to our diagram configuration $s$ doesn't lie on the topmost
row, so there was a black $b'$ square above $s$ previously. Since $b'$
lies on the rightmost column it must have been paired up with the square
below it, $s$, but this means that $s$ doesn't exist a
contradiction. \vspace{0.2cm}

Now suppose $s$ does lie on the bottom row. Due to our diagram configuration
there is always a black square $b$ to the right of $s$. The only issue
occurs when this black square $b$ happens to be the one removed from the
chessboard. If this is the case then there was a black square above $s$
call it $b'$. At that point in time the square below $b'$ was available to
take and since our algorithm first prioritises taking the square below
$b'$ we would have removed $b'$ and $s$ previously. But this means that
$s$ doesn't exist, a contradiction. $\square$
\\ \bigskip

(b) After experimenting we find that it seems to be that only the
very middle square can be left uncovered. Let's prove our guess.
\vspace{0.2cm}

Suppose FTSOC that the middle square has been covered, by (WLOG) a
horizontal $3\times 1$ triomino. The following claim will be useful to us.
\vspace{0.2cm}

\textit{Claim 1: For any row, at most one horizontal $3\times 1$ triomino
can be placed there. Similarly for any column, at most one vertical
$1\times 3$ triomino can be placed there.}
\vspace{0.2cm}

\textit{Proof:} We will show the first statement to be true, this also
gives the second statement to be true since we can rotate the board by
$90^\circ$. In a row there are 5 squares. Each $3\times 1$ triomino uses
exactly 3 squares. If we could place $t \geq 2$ triominos on a row then
we require
$
3t \leq 5
$
but since $t \geq 2$ this inequality is false. Thus $t \leq 1$ as desired.
$\square$ \\

The horizontal triomino that we have placed to cover the middle square also
splits precisely 3 columns into $1\times 2$ rectangles that are above and
below of our triomino. The diagram below illustrates the idea.

\begin{center}
\begin{asy}
/* Geogebra to Asymptote conversion, documentation at artofproblemsolving.com/Wiki go to User:Azjps/geogebra */
import graph; size(150);
real labelscalefactor = 0.5; /* changes label-to-point distance */
pen dps = linewidth(0.7) + fontsize(10); defaultpen(dps); /* default pen style */
pen dotstyle = black; /* point style */
real xmin = 0.2424095295677375, xmax = 6.664363890248977, ymin = 0.4060647819051748, ymax = 6.327145985906095;  /* image dimensions */
pen qqzzff = rgb(0.,0.6,1.);

filldraw((1.,6.)--(2.,6.)--(2.,5.)--(1.,5.)--cycle, black + opacity(1.0));
filldraw((3.,6.)--(4.,6.)--(4.,5.)--(3.,5.)--cycle, black + opacity(1.0));
filldraw((5.,6.)--(6.,6.)--(6.,5.)--(5.,5.)--cycle, black + opacity(1.0));
filldraw((2.,5.)--(3.,5.)--(3.,4.)--(2.,4.)--cycle, black + opacity(1.0));
filldraw((4.,5.)--(5.,5.)--(5.,4.)--(4.,4.)--cycle, black + opacity(1.0));
filldraw((5.,4.)--(6.,4.)--(6.,3.)--(5.,3.)--cycle, black + opacity(1.0));
filldraw((1.,4.)--(2.,4.)--(2.,3.)--(1.,3.)--cycle, black + opacity(1.0));
filldraw((3.,4.)--(4.,4.)--(4.,3.)--(3.,3.)--cycle, black + opacity(1.0));
filldraw((4.,3.)--(5.,3.)--(5.,2.)--(4.,2.)--cycle, black + opacity(1.0));
filldraw((2.,3.)--(3.,3.)--(3.,2.)--(2.,2.)--cycle, black + opacity(1.0));
filldraw((5.,2.)--(6.,2.)--(6.,1.)--(5.,1.)--cycle, black + opacity(1.0));
filldraw((3.,2.)--(4.,2.)--(4.,1.)--(3.,1.)--cycle, black + opacity(1.0));
filldraw((1.,2.)--(2.,2.)--(2.,1.)--(1.,1.)--cycle, black + opacity(1.0));
filldraw(circle((2.5,3.5), 0.25), qqzzff + opacity(1.0),  qqzzff);
filldraw(circle((3.5,3.5), 0.25), qqzzff + opacity(1.0),  qqzzff);
filldraw(circle((4.5,3.5), 0.25), qqzzff + opacity(1.0),  qqzzff);
/* draw figures */
draw((1.,6.)--(2.,6.));
draw((2.,6.)--(2.,5.), linewidth(3.0) + red);
draw((2.,5.)--(1.,5.));
draw((1.,5.)--(1.,6.));
draw((3.,6.)--(4.,6.));
draw((4.,6.)--(4.,5.));
draw((4.,5.)--(3.,5.));
draw((3.,5.)--(3.,6.), linewidth(3.0) + red);
draw((5.,6.)--(6.,6.));
draw((6.,6.)--(6.,5.));
draw((6.,5.)--(5.,5.));
draw((5.,5.)--(5.,6.));
draw((2.,5.)--(3.,5.));
draw((3.,5.)--(3.,4.), linewidth(3.0) + red);
draw((3.,4.)--(2.,4.), linewidth(3.0) + red);
draw((2.,4.)--(2.,5.), linewidth(3.0) + red);
draw((4.,5.)--(5.,5.));
draw((5.,5.)--(5.,4.));
draw((5.,4.)--(4.,4.));
draw((4.,4.)--(4.,5.));
draw((5.,4.)--(6.,4.));
draw((6.,4.)--(6.,3.));
draw((6.,3.)--(5.,3.));
draw((5.,3.)--(5.,4.));
draw((1.,4.)--(2.,4.));
draw((2.,4.)--(2.,3.));
draw((2.,3.)--(1.,3.));
draw((1.,3.)--(1.,4.));
draw((3.,4.)--(4.,4.));
draw((4.,4.)--(4.,3.));
draw((4.,3.)--(3.,3.));
draw((3.,3.)--(3.,4.));
draw((4.,3.)--(5.,3.));
draw((5.,3.)--(5.,2.));
draw((5.,2.)--(4.,2.));
draw((4.,2.)--(4.,3.));
draw((2.,3.)--(3.,3.), linewidth(3.0) + red);
draw((3.,3.)--(3.,2.), linewidth(3.0) + red);
draw((3.,2.)--(2.,2.));
draw((2.,2.)--(2.,3.), linewidth(3.0) + red);
draw((5.,2.)--(6.,2.));
draw((6.,2.)--(6.,1.));
draw((6.,1.)--(5.,1.));
draw((5.,1.)--(5.,2.));
draw((3.,2.)--(4.,2.));
draw((4.,2.)--(4.,1.));
draw((4.,1.)--(3.,1.));
draw((3.,1.)--(3.,2.), linewidth(3.0) + red);
draw((1.,2.)--(2.,2.));
draw((2.,2.)--(2.,1.), linewidth(3.0) + red);
draw((2.,1.)--(1.,1.));
draw((1.,1.)--(1.,2.));
draw((6.,6.)--(6.,1.));
draw((1.,1.)--(1.,6.));
draw((2.,1.)--(3.,1.), linewidth(3.0) + red);
draw((2.,6.)--(3.,6.), linewidth(3.0) + red);
draw((4.,6.)--(5.,6.));
/* dots and labels */
clip((xmin,ymin)--(xmin,ymax)--(xmax,ymax)--(xmax,ymin)--cycle);
/* end of picture */
\end{asy}
\end{center}

This means for those three columns, there is no longer any way to place a
vertical $1\times 3$ triomino. Now using claim 1,
on the entire board we can have at most 5 horizontal triominos, and at most
$5 - 3 = 2$ vertical triominos. In total we can place at most 7 triominos but we
wanted to place 8 triominos, a contradiction. \vspace{0.2cm}

Now we show a way to place the 8 triominoes with just the very middle
uncovered. $\square$

\begin{center}
\begin{asy}
/* Geogebra to Asymptote conversion, documentation at artofproblemsolving.com/Wiki go to User:Azjps/geogebra */
import graph; size(150);
real labelscalefactor = 0.5; /* changes label-to-point distance */
pen dps = linewidth(0.7) + fontsize(10); defaultpen(dps); /* default pen style */
pen dotstyle = black; /* point style */
real xmin = 0.4427549836330455, xmax = 6.556224887729826, ymin = 0.5053241563496622, ymax = 6.446328388249316;  /* image dimensions */
pen qqzzff = rgb(0.,0.6,1.); pen qqzzqq = rgb(0.,0.6,0.); pen ffdxqq = rgb(1.,0.8431372549019608,0.); pen ttzzqq = rgb(0.2,0.6,0.);

filldraw((1.,6.)--(2.,6.)--(2.,5.)--(1.,5.)--cycle, black + opacity(1.0));
filldraw((3.,6.)--(4.,6.)--(4.,5.)--(3.,5.)--cycle, black + opacity(1.0));
filldraw((5.,6.)--(6.,6.)--(6.,5.)--(5.,5.)--cycle, black + opacity(1.0));
filldraw((2.,5.)--(3.,5.)--(3.,4.)--(2.,4.)--cycle, black + opacity(1.0));
filldraw((4.,5.)--(5.,5.)--(5.,4.)--(4.,4.)--cycle, black + opacity(1.0));
filldraw((5.,4.)--(6.,4.)--(6.,3.)--(5.,3.)--cycle, black + opacity(1.0));
filldraw((1.,4.)--(2.,4.)--(2.,3.)--(1.,3.)--cycle, black + opacity(1.0));
filldraw((3.,4.)--(4.,4.)--(4.,3.)--(3.,3.)--cycle, black + opacity(1.0));
filldraw((4.,3.)--(5.,3.)--(5.,2.)--(4.,2.)--cycle, black + opacity(1.0));
filldraw((2.,3.)--(3.,3.)--(3.,2.)--(2.,2.)--cycle, black + opacity(1.0));
filldraw((5.,2.)--(6.,2.)--(6.,1.)--(5.,1.)--cycle, black + opacity(1.0));
filldraw((3.,2.)--(4.,2.)--(4.,1.)--(3.,1.)--cycle, black + opacity(1.0));
filldraw((1.,2.)--(2.,2.)--(2.,1.)--(1.,1.)--cycle, black + opacity(1.0));
filldraw(circle((2.5,1.5), 0.25), red + opacity(1.0),  red);
filldraw(circle((3.5,1.5), 0.25), red + opacity(1.0),  red);
filldraw(circle((1.5,1.5), 0.25), red + opacity(1.0),  red);
filldraw(circle((1.5,2.5), 0.25), qqzzff + opacity(1.0),  qqzzff);
filldraw(circle((2.5,2.5), 0.25), qqzzff + opacity(1.0),  qqzzff);
filldraw(circle((3.5,2.5), 0.25), qqzzff + opacity(1.0),  qqzzff);
filldraw(circle((1.5,3.5), 0.25), qqzzqq + opacity(1.0),  qqzzqq);
filldraw(circle((1.5,4.5), 0.25), qqzzqq + opacity(1.0),  qqzzqq);
filldraw(circle((1.5,5.5), 0.25), qqzzqq + opacity(1.0),  qqzzqq);
filldraw(circle((2.5,3.5), 0.25), ffdxqq + opacity(1.0),  ffdxqq);
filldraw(circle((2.5,4.5), 0.25), ffdxqq + opacity(1.0),  ffdxqq);
filldraw(circle((2.5,5.5), 0.25), ffdxqq + opacity(1.0),  ffdxqq);
filldraw(circle((3.5,4.5), 0.25), red + opacity(1.0),  red);
filldraw(circle((4.5,4.5), 0.25), red + opacity(1.0),  red);
filldraw(circle((5.5,4.5), 0.25), red + opacity(1.0),  red);
filldraw(circle((3.5,5.5), 0.25), qqzzff + opacity(1.0),  qqzzff);
filldraw(circle((4.5,5.5), 0.25), qqzzff + opacity(1.0),  qqzzff);
filldraw(circle((5.5,5.5), 0.25), qqzzff + opacity(1.0),  qqzzff);
filldraw(circle((4.5,3.5), 0.25), ttzzqq + opacity(1.0),  ttzzqq);
filldraw(circle((4.5,2.5), 0.25), ttzzqq + opacity(1.0),  ttzzqq);
filldraw(circle((4.5,1.5), 0.25), ttzzqq + opacity(1.0),  ttzzqq);
filldraw(circle((5.5,3.5), 0.25), ffdxqq + opacity(1.0),  ffdxqq);
filldraw(circle((5.5,2.5), 0.25), ffdxqq + opacity(1.0),  ffdxqq);
filldraw(circle((5.5,1.5), 0.25), ffdxqq + opacity(1.0),  ffdxqq);
/* draw figures */
draw((1.,6.)--(2.,6.));
draw((2.,6.)--(2.,5.));
draw((2.,5.)--(1.,5.));
draw((1.,5.)--(1.,6.));
draw((3.,6.)--(4.,6.));
draw((4.,6.)--(4.,5.));
draw((4.,5.)--(3.,5.));
draw((3.,5.)--(3.,6.));
draw((5.,6.)--(6.,6.));
draw((6.,6.)--(6.,5.));
draw((6.,5.)--(5.,5.));
draw((5.,5.)--(5.,6.));
draw((2.,5.)--(3.,5.));
draw((3.,5.)--(3.,4.));
draw((3.,4.)--(2.,4.));
draw((2.,4.)--(2.,5.));
draw((4.,5.)--(5.,5.));
draw((5.,5.)--(5.,4.));
draw((5.,4.)--(4.,4.));
draw((4.,4.)--(4.,5.));
draw((5.,4.)--(6.,4.));
draw((6.,4.)--(6.,3.));
draw((6.,3.)--(5.,3.));
draw((5.,3.)--(5.,4.));
draw((1.,4.)--(2.,4.));
draw((2.,4.)--(2.,3.));
draw((2.,3.)--(1.,3.));
draw((1.,3.)--(1.,4.));
draw((3.,4.)--(4.,4.));
draw((4.,4.)--(4.,3.));
draw((4.,3.)--(3.,3.));
draw((3.,3.)--(3.,4.));
draw((4.,3.)--(5.,3.));
draw((5.,3.)--(5.,2.));
draw((5.,2.)--(4.,2.));
draw((4.,2.)--(4.,3.));
draw((2.,3.)--(3.,3.));
draw((3.,3.)--(3.,2.));
draw((3.,2.)--(2.,2.));
draw((2.,2.)--(2.,3.));
draw((5.,2.)--(6.,2.));
draw((6.,2.)--(6.,1.));
draw((6.,1.)--(5.,1.));
draw((5.,1.)--(5.,2.));
draw((3.,2.)--(4.,2.));
draw((4.,2.)--(4.,1.));
draw((4.,1.)--(3.,1.));
draw((3.,1.)--(3.,2.));
draw((1.,2.)--(2.,2.));
draw((2.,2.)--(2.,1.));
draw((2.,1.)--(1.,1.));
draw((1.,1.)--(1.,2.));
draw((6.,6.)--(6.,1.));
draw((1.,1.)--(1.,6.));
draw((2.,1.)--(3.,1.));
draw((2.,6.)--(3.,6.));
draw((4.,6.)--(5.,6.));
draw((4.,1.)--(5.,1.));
/* dots and labels */
clip((xmin,ymin)--(xmin,ymax)--(xmax,ymax)--(xmax,ymin)--cycle);
/* end of picture */
\end{asy}
\end{center}

\bigskip
(c) Again after experimenting we find that only the corners of the board
can be left uncovered in this case. Let's prove our guess.
\vspace{0.2cm}

Similar to claim 1 (b), we can place at most 5 horizontal (or vertical)
$4 \times 1$ tetrominoes on the board, because each row (or column)
only has enough squares to accomodate for at most one
$4 \times 1$ tetromino. Suppose we use $v$ vertical tetrominoes and
$h$ horizontal tetrominoes. At least one of $v,h$ has to be $\geq 3$
because otherwise
$$
6 = v + h < 3 + 3 = 6
$$
is a contradiction. Suppose WLOG that $v \geq 3$.
\vspace{0.2cm}

Each of those vertical pieces always covers the middle three squares of
that column. A picture helps to illustrate.

\begin{center}
\begin{asy}
/* Geogebra to Asymptote conversion, documentation at artofproblemsolving.com/Wiki go to User:Azjps/geogebra */
import graph; size(150);
real labelscalefactor = 0.5; /* changes label-to-point distance */
pen dps = linewidth(0.7) + fontsize(10); defaultpen(dps); /* default pen style */
pen dotstyle = black; /* point style */
real xmin = 0.5180791366329588, xmax = 6.550211502796489, ymin = 0.5291618275609714, ymax = 6.539144568075509;  /* image dimensions */
pen ttttff = rgb(0.2,0.2,1.);

filldraw((1.,6.)--(2.,6.)--(2.,5.)--(1.,5.)--cycle, black + opacity(1.0));
filldraw((3.,6.)--(4.,6.)--(4.,5.)--(3.,5.)--cycle, black + opacity(1.0));
filldraw((5.,6.)--(6.,6.)--(6.,5.)--(5.,5.)--cycle, black + opacity(1.0));
filldraw((2.,5.)--(3.,5.)--(3.,4.)--(2.,4.)--cycle, black + opacity(1.0));
filldraw((4.,5.)--(5.,5.)--(5.,4.)--(4.,4.)--cycle, black + opacity(1.0));
filldraw((5.,4.)--(6.,4.)--(6.,3.)--(5.,3.)--cycle, black + opacity(1.0));
filldraw((3.,4.)--(4.,4.)--(4.,3.)--(3.,3.)--cycle, black + opacity(1.0));
filldraw((2.,4.)--(2.,3.)--(1.,3.)--(1.,4.)--cycle, black + opacity(1.0));
filldraw((4.,3.)--(5.,3.)--(5.,2.)--(4.,2.)--cycle, black + opacity(1.0));
filldraw((2.,3.)--(3.,3.)--(3.,2.)--(2.,2.)--cycle, black + opacity(1.0));
filldraw((1.,2.)--(2.,2.)--(2.,1.)--(1.,1.)--cycle, black + opacity(1.0));
filldraw((3.,2.)--(4.,2.)--(4.,1.)--(3.,1.)--cycle, black + opacity(1.0));
filldraw((5.,2.)--(6.,2.)--(6.,1.)--(5.,1.)--cycle, black + opacity(1.0));
filldraw(circle((2.5,1.5), 0.25), ttttff + opacity(1.0),  ttttff);
filldraw(circle((2.5,2.5), 0.25), ttttff + opacity(1.0),  ttttff);
filldraw(circle((2.5,3.5), 0.25), ttttff + opacity(1.0),  ttttff);
filldraw(circle((2.5,4.5), 0.25), ttttff + opacity(1.0),  ttttff);
filldraw(circle((4.5,2.5), 0.25), ttttff + opacity(1.0),  ttttff);
filldraw(circle((4.5,3.5), 0.25), ttttff + opacity(1.0),  ttttff);
filldraw(circle((4.5,4.5), 0.25), ttttff + opacity(1.0),  ttttff);
filldraw(circle((4.5,5.5), 0.25), ttttff + opacity(1.0),  ttttff);
filldraw(circle((5.5,2.5), 0.25), ttttff + opacity(1.0),  ttttff);
filldraw(circle((5.5,3.5), 0.25), ttttff + opacity(1.0),  ttttff);
filldraw(circle((5.5,4.5), 0.25), ttttff + opacity(1.0),  ttttff);
filldraw(circle((5.5,5.5), 0.25), ttttff + opacity(1.0),  ttttff);
/* draw figures */
draw((1.,6.)--(2.,6.));
draw((2.,6.)--(2.,5.));
draw((2.,5.)--(1.,5.));
draw((1.,5.)--(1.,6.));
draw((3.,6.)--(4.,6.));
draw((4.,6.)--(4.,5.));
draw((4.,5.)--(3.,5.));
draw((3.,5.)--(3.,6.));
draw((5.,6.)--(6.,6.));
draw((6.,6.)--(6.,5.));
draw((6.,5.)--(5.,5.), linewidth(2.8) + red);
draw((5.,5.)--(5.,6.));
draw((2.,5.)--(3.,5.), linewidth(2.8) + red);
draw((3.,5.)--(3.,4.), linewidth(2.8) + red);
draw((3.,4.)--(2.,4.));
draw((2.,4.)--(2.,5.), linewidth(2.8) + red);
draw((4.,5.)--(5.,5.), linewidth(2.8) + red);
draw((5.,5.)--(5.,4.), linewidth(2.8) + red);
draw((5.,4.)--(4.,4.));
draw((4.,4.)--(4.,5.), linewidth(2.8) + red);
draw((5.,4.)--(6.,4.));
draw((6.,4.)--(6.,3.), linewidth(2.8) + red);
draw((6.,3.)--(5.,3.));
draw((5.,3.)--(5.,4.), linewidth(2.8) + red);
draw((3.,4.)--(4.,4.));
draw((4.,4.)--(4.,3.), linewidth(2.8) + red);
draw((4.,3.)--(3.,3.));
draw((3.,3.)--(3.,4.), linewidth(2.8) + red);
draw((2.,4.)--(2.,3.), linewidth(2.8) + red);
draw((2.,3.)--(1.,3.));
draw((1.,3.)--(1.,4.));
draw((1.,4.)--(2.,4.));
draw((4.,3.)--(5.,3.));
draw((5.,3.)--(5.,2.), linewidth(2.8) + red);
draw((5.,2.)--(4.,2.), linewidth(2.8) + red);
draw((4.,2.)--(4.,3.), linewidth(2.8) + red);
draw((2.,3.)--(3.,3.));
draw((3.,3.)--(3.,2.), linewidth(2.8) + red);
draw((3.,2.)--(2.,2.), linewidth(2.8) + red);
draw((2.,2.)--(2.,3.), linewidth(2.8) + red);
draw((1.,2.)--(2.,2.));
draw((2.,2.)--(2.,1.));
draw((2.,1.)--(1.,1.));
draw((1.,1.)--(1.,2.));
draw((3.,2.)--(4.,2.));
draw((4.,2.)--(4.,1.));
draw((4.,1.)--(3.,1.));
draw((3.,1.)--(3.,2.));
draw((5.,2.)--(6.,2.), linewidth(2.8) + red);
draw((6.,2.)--(6.,1.));
draw((6.,1.)--(5.,1.));
draw((5.,1.)--(5.,2.));
draw((2.,6.)--(3.,6.));
draw((4.,6.)--(5.,6.));
draw((6.,5.)--(6.,4.), linewidth(2.8) + red);
draw((6.,3.)--(6.,2.), linewidth(2.8) + red);
draw((5.,1.)--(4.,1.));
draw((3.,1.)--(2.,1.));
draw((1.,2.)--(1.,3.));
draw((1.,4.)--(1.,5.));
/* dots and labels */
clip((xmin,ymin)--(xmin,ymax)--(xmax,ymax)--(xmax,ymin)--cycle);
/* end of picture */
\end{asy}
\end{center}

So now we can no longer place horizontal tetrominos along those middle three
rows because there are only two squares left uncovered whilst we need
at least 4 to place our tetrominos. Thus we must place horizontal tetrominos
along the top or the bottom of the board. \vspace{0.2cm}

Actually we can place at most one of these horizontal tetrominos on the
board since if we place two then we will be unable to place any vertical
tetrominoes on the central three columns, by a similar arguement to the one
above. Thus in order to place 6 tetrominos we must use 5 vertical
tetrominoes and 1 horizontal tetromino. \vspace{-0.2cm}

There must be a vertical tetromino in every column, since we have a
$5\times 5$ board. Once we place the horizontal tetromino along the top or
the bottom, in each of the columns that the horizontal tetromino passes
through, there is now only one way to place the vertical tetromino for that
column. The only degree of freedom we have is where we place the last
vertical tetromino.

\begin{center}
\begin{asy}
/* Geogebra to Asymptote conversion, documentation at artofproblemsolving.com/Wiki go to User:Azjps/geogebra */
import graph; size(150);
real labelscalefactor = 0.5; /* changes label-to-point distance */
pen dps = linewidth(0.7) + fontsize(10); defaultpen(dps); /* default pen style */
pen dotstyle = black; /* point style */
real xmin = 0.5029768588989332, xmax = 6.504391120413052, ymin = 0.5071424296121633, ymax = 6.500797522161037;  /* image dimensions */
pen ttttff = rgb(0.2,0.2,1.); pen ffdxqq = rgb(1.,0.8431372549019608,0.); pen bubtba = rgb(0.7058823529411765,0.7019607843137254,0.7294117647058823); pen rcrcrf = rgb(0.10980392156862745,0.10980392156862745,0.12156862745098039);

filldraw((1.,6.)--(2.,6.)--(2.,5.)--(1.,5.)--cycle, black);
filldraw((3.,6.)--(4.,6.)--(4.,5.)--(3.,5.)--cycle, black);
filldraw((5.,6.)--(6.,6.)--(6.,5.)--(5.,5.)--cycle, black);
filldraw((2.,5.)--(3.,5.)--(3.,4.)--(2.,4.)--cycle, black);
filldraw((4.,5.)--(5.,5.)--(5.,4.)--(4.,4.)--cycle, black);
filldraw((5.,4.)--(6.,4.)--(6.,3.)--(5.,3.)--cycle, black);
filldraw((3.,4.)--(4.,4.)--(4.,3.)--(3.,3.)--cycle, black);
filldraw((2.,4.)--(2.,3.)--(1.,3.)--(1.,4.)--cycle, black);
filldraw((4.,3.)--(5.,3.)--(5.,2.)--(4.,2.)--cycle, black);
filldraw((2.,3.)--(3.,3.)--(3.,2.)--(2.,2.)--cycle, black);
filldraw((1.,2.)--(2.,2.)--(2.,1.)--(1.,1.)--cycle, black);
filldraw((3.,2.)--(4.,2.)--(4.,1.)--(3.,1.)--cycle, black);
filldraw((5.,2.)--(6.,2.)--(6.,1.)--(5.,1.)--cycle, black);
filldraw(circle((4.5,2.5), 0.25), ttttff,  ttttff);
filldraw(circle((4.5,3.5), 0.25), ttttff,  ttttff);
filldraw(circle((4.5,4.5), 0.25), ttttff,  ttttff);
filldraw(circle((4.5,5.5), 0.25), ttttff,  ttttff);
filldraw(circle((3.5,2.5), 0.25), red,  red);
filldraw(circle((3.5,3.5), 0.25), red,  red);
filldraw(circle((3.5,4.5), 0.25), red,  red);
filldraw(circle((3.5,5.5), 0.25), red,  red);
filldraw(circle((2.5,2.5), 0.25), ttttff,  ttttff);
filldraw(circle((2.5,3.5), 0.25), ttttff,  ttttff);
filldraw(circle((2.5,4.5), 0.25), ttttff,  ttttff);
filldraw(circle((2.5,5.5), 0.25), ttttff,  ttttff);
filldraw(circle((1.5,2.5), 0.25), red,  red);
filldraw(circle((1.5,3.5), 0.25), red,  red);
filldraw(circle((1.5,4.5), 0.25), red,  red);
filldraw(circle((1.5,5.5), 0.25), red,  red);
filldraw(circle((1.5,1.5), 0.25), ffdxqq,  ffdxqq);
filldraw(circle((2.5,1.5), 0.25), ffdxqq,  ffdxqq);
filldraw(circle((3.5,1.5), 0.25), ffdxqq,  ffdxqq);
filldraw(circle((4.5,1.5), 0.25), ffdxqq,  ffdxqq);

/* draw figures */
draw((1.,6.)--(2.,6.));
draw((2.,6.)--(2.,5.));
draw((2.,5.)--(1.,5.));
draw((1.,5.)--(1.,6.));
draw((3.,6.)--(4.,6.));
draw((4.,6.)--(4.,5.));
draw((4.,5.)--(3.,5.));
draw((3.,5.)--(3.,6.));
draw((5.,6.)--(6.,6.));
draw((6.,6.)--(6.,5.));
draw((6.,5.)--(5.,5.));
draw((5.,5.)--(5.,6.));
draw((2.,5.)--(3.,5.));
draw((3.,5.)--(3.,4.));
draw((3.,4.)--(2.,4.));
draw((2.,4.)--(2.,5.));
draw((4.,5.)--(5.,5.));
draw((5.,5.)--(5.,4.));
draw((5.,4.)--(4.,4.));
draw((4.,4.)--(4.,5.));
draw((5.,4.)--(6.,4.));
draw((6.,4.)--(6.,3.));
draw((6.,3.)--(5.,3.));
draw((5.,3.)--(5.,4.));
draw((3.,4.)--(4.,4.));
draw((4.,4.)--(4.,3.));
draw((4.,3.)--(3.,3.));
draw((3.,3.)--(3.,4.));
draw((2.,4.)--(2.,3.));
draw((2.,3.)--(1.,3.));
draw((1.,3.)--(1.,4.));
draw((1.,4.)--(2.,4.));
draw((4.,3.)--(5.,3.));
draw((5.,3.)--(5.,2.));
draw((5.,2.)--(4.,2.));
draw((4.,2.)--(4.,3.));
draw((2.,3.)--(3.,3.));
draw((3.,3.)--(3.,2.));
draw((3.,2.)--(2.,2.));
draw((2.,2.)--(2.,3.));
draw((1.,2.)--(2.,2.));
draw((2.,2.)--(2.,1.));
draw((2.,1.)--(1.,1.));
draw((1.,1.)--(1.,2.));
draw((3.,2.)--(4.,2.));
draw((4.,2.)--(4.,1.));
draw((4.,1.)--(3.,1.));
draw((3.,1.)--(3.,2.));
draw((5.,2.)--(6.,2.));
draw((6.,2.)--(6.,1.));
draw((6.,1.)--(5.,1.));
draw((5.,1.)--(5.,2.));
draw((2.,6.)--(3.,6.));
draw((4.,6.)--(5.,6.));
draw((6.,5.)--(6.,4.));
draw((6.,3.)--(6.,2.));
draw((5.,1.)--(4.,1.));
draw((3.,1.)--(2.,1.));
draw((1.,2.)--(1.,3.));
draw((1.,4.)--(1.,5.));
/* dots and labels */
clip((xmin,ymin)--(xmin,ymax)--(xmax,ymax)--(xmax,ymin)--cycle);
/* end of picture */
\end{asy}
\end{center}

in any case wherever we place the last vertical tetromino the uncovered
square will be in the corner. Finally it is easy to demonstrate by
construction that every corner can in fact be left uncovered (I omit
diagrams here because I'm tired of dealing with Geogebra by now. :p).
$\blacksquare$
\\

\textit{Remarks: Interstingly enough, for me at least, part (c) was the
easiest, then part (b), then part (a) was the hardest, whilst usually it's
the reverse way round in terms of difficulties.}
\\

\hline
\vspace{0.2cm}

\textit{Addendums (after having read solutions):} This is a nice game-theory
type problem. Each of my three solutions look really different to the
official solutions so let's discuss them one by one. \vspace{0.2cm}

For part (a), I remembered the last problems we did one chessboards like
this and that it should be possible to just remove any black square from
the board and then tile the subsequent board with $2\times 1$ dominoes.
Playing around with the problem I found that there are many ways to do these
tilings so I just guessed that a greedy solution might work. Then finally
divising the algorithm and proving it's correctness feels like a standard
codeforces problem. The case bashing in the official solutions is probably
a much clearer solution. \vspace{0.2cm}

For parts (b) and (c) my two solutions are pretty similar, somehow some
configuration must be forced then we case bash to show the final result.
The colouring certain tiles solution from the official sols here are
particularly elegant.
\end{solution}
\newpage

\begin{problem}{Problem 7}
Let us call a positive integer $n$ "mischievous" if both $n$ and $n+1$ are
divisible by perfect cubes greater than 1. \vspace{0.2cm}

Show that at least 1\% of positive integers are mischievous.
\vspace{0.2cm}

\textit{
You might wonder how you can talk about 1\% of an infinite number of
numbers. What is meant is that once you take sufficiently many numbers,
at least 1\% are mischievous. Formally, you want to show that there is
some $N$ such that any set of $N$ consectutive numbers, then 1\% of
them are mischievous.
}
\end{problem}

\begin{hint}{Initial thoughts and motivations}
\textit{Motivations.}
\textit{
Let $p,q,n$ be positive integers. Employing
some wishful thinking it would be nice if I could show that there
is a family of solutions to}
$$
p^3 - nq^3 = 1
$$
\textit{
because then I could set $n+1 = nq^3 + 1, n = nq^3$ and this way both of
them are divisible by perfect cubes. Looking at this equation it looks
similar to the usual Pell equation form but instead of squares we have
cubes.} \vspace{0.2cm}

\textit{Next I naturally thought about the identity}
$$
n^3 + 1 = (n+1)(n^2 - n + 1)
$$
\textit{
since it looks similar to what I wanted above. Actually if I set
$n+1 = t^3$ for some $t \in \mathbb{N}$ then this would be exactly what
I wanted and I believe it actually works. Let's have a go.} \\

Let $n \in \mathbb{N}$. Consider the following fact
$$
n^3 + 1 = (n + 1)(n^2 - n + 1)
$$
Let $k \in \mathbb{N}$. Set $n = k^3 - 1$. Now we have
$$
(k^3 - 1)^3 + 1 = k^3 \left[ (k^3 - 1)^2 - (k^3 - 1) + 1 \right].
$$
Thus since $(k^3 - 1)^3 + 1$ is divisible by a perfect cube, the number
$(k^3 - 1)^3$ is mischievous. \\

\textit{
Unfortunately although this shows that there are infinitely many
mischievous numbers, the gap between two consecutive numbers of this form
shoots off to infinity for large enough $k$, so this approach won't work
to show anything about the density of mischievous numbers. Circling back
to the pell equation approach, we could consider the equation
}
$$
x^6 - n y^6 = 1
$$
\textit{but analysing the density of these solutions seems more complicated
than I'd hope, so I think it's better to look for something else.
If I can factor the number $kn^3 + 1$ into (cube)(something), then
this would be desirable. I wondered what if for that (cube) I just try and
force $kn^3 + 1$ to be divisible by some small cube say 8. If you go down
this path you can show, for instance, that the numbers
$7^3, 15^3, 23^3, \ldots$ are
all mischevious. These numbers are denser than the last family,
this seems like progress. Now what if I replaced 8 with 27. Then you can
show that the numbers $26^3, 53^3, 80^3, \ldots$ are all mischevious.
Let's try and generalise. I conjecture that all numbers of the form
}
$$
(k^3 - 1 + tk)^3, \quad (t = 0,1,2,\ldots)
$$
\textit{
are all mischievous. Still though I think these numbers aren't dense enough.
An equivalent formulation of my idea is to study the modular equation
}
\[
\begin{aligned}
kn^3 + 1 &\equiv 0 \ (\text{mod } p^3) \\
kn^3 &\equiv -1 \ (\text{mod } p^3)
\end{aligned}
\]
\textit{Let's restrict our analysis to say $p=2$. It might be helpful to
employ Euler's theorem here, so we compute $\varphi(8) = 8(1 - 1/2) = 4$.
Then for $\gcd(n, 8) = 1$ we can write,}
\[
\begin{aligned}
kn^3 &\equiv -1 \ (\text{mod } 8) \\
kn^4 &\equiv -n \ (\text{mod } 8) \\
k &\equiv -n \ (\text{mod } 8) \\
\end{aligned}
\]
\textit{So actually pairs of the form $(k, n) = (-n, n) \ \text{mod } 8$,
$\gcd(n,8) = 1$ should always give us mischevious numbers. Numbers of this
class still aren't dense enough so maybe lets try to generalise.
We note that $\varphi(p^3) = p^3(1-1/p) = p^2(p-1)$.
}
\[
\begin{aligned}
kn^3 &\equiv -1 \ (\text{mod } p^3) \\
kn^3\cdot n^{p^3 - p^2 - 3} &\equiv -n^{p^3 - p^2 - 3} \ (\text{mod } p^3)\\
k &\equiv -n^{p^3 - p^2 - 3} \ (\text{mod } p^3)\\
\end{aligned}
\]
\textit{Thus pairs of the form}
$$
(k,n) = (-n^{p^3 - p^2 - 3}, n) \ \text{mod } p^3
$$
\textit{always produce mischievous numbers for $\gcd(p,n)=1$. Modulo
$p^3$ the number $kn^3 + 1$ then looks like}
$$
-n^{p^3 - p^2} + 1 \equiv n^{p^3 - p^2} - 1 \equiv 0
$$
\textit{and we can make the statement that numbers of the form}
$$
n^{p^3 - p^2} - 1, \quad (n = 1,2,\ldots,p-1)
$$
\textit{are always mischievous. There is the corollary that
for instance numbers of the form $n^4 - 1$ are also always mischievous,
for $\gcd(n, 8) = 1$. Roughly we can say that the gap between two
consecutive numbers of the form $n^4 - 1$ grows like $\mathcal{O}(n^3)$.
More generally the gap between two consecutive numbers of the form
$n^{p^3 - p^2} - 1$, grows like $\mathcal{O}(n^{p^3 - p^2 -1})$. If I can
improve this gap to even $\mathcal{O}(\text{some large constant})$
this would be great. Actually I see a way to improve this gap to
$\mathcal{O}(n^2)$. Here is the motivation.} \vspace{0.2cm}

\textit{
I seek to find some number $n = kp^3 - 1$ with the condition
$v_q (kp^3 - 1) \geq 3$ for some prime $q$. If swap $k$ with $k^3$ I
could envoke LTE here. In particular if I let $k, p$ be odd the following
happens via LTE}
$$
\nu_2\left( k^3 p^3 - 1 \right) = \nu_2 (kp - 1)
$$
\textit{Now if I let $kp - 1 = 8x$ this guarentees $\nu_2(kp - 1) \geq 3$
and thus numbers of the form $(8x + 1)^3 - 1$ are mischevious. Actually
re-reading over this that statement is quite obviously true since
reducing mod 8 gives $1 - 1 = 0$. Moreover it is indeed the case that the
gap between two consecutive numbers of this form grows like
$\mathcal{O}(x^2)$ which is an improvement! We can generalise this idea
to get that numbers of the form $(p^3 x  + 1)^3 - 1$ are mischievous also,
where $p$ is a fixed prime.
}
\vspace{0.2cm}

\textit{From here it is natural to try and improve the gap to linear
or constant. Another idea is to study the solutions to }
$$
27x + 8y = 1.
$$
\textit{We can find solutions by say using Bézout's theorem. After
playing around with the equation we notice the following pattern}
\[
\begin{aligned}
27(3) - 8(10) &= 1 \\
27(3) - 8(10) - 27(8) + 27(8) &= 1\\
27(11) - 8(37) &= 1 \\
27(11) - 8(37) - 27(8) + 27(8) &= 1 \\
27(19) - 8(64) &= 1 \\
\cdots
\end{aligned}
\]
\textit{and so we find that the numbers $8(10), 8(37), 8(64), \ldots$
are all mischievous, since adding 1 gives a number divisible by 27.
Now this is very promising since these numbers are
much more dense than the numbers we were considering before. If this
conjecture is true than we will get that all numbers of the form
$8(10 + 27n)$ are mischievous. Now the gap between two consecutive numbers
of this form is $27(8) = 216$, so I could use this to show that
$1/216 \approx 0.46 \%$ of numbers are mischievous.
We can use similar
methods to show that all numbers of the form $27(37 + 125n)$ are
mischievous, but this method seem to be giving us diminishing returns as
it only gives that $1/3375 \approx 0.03 \%$ of numbers are
mischievous. Similarly
all numbers of the form $125(3 + 8n)$ are mischievous which gives
$1/1000 = 0.1\%$ of numbers being mischevious.
Indeed if I just spam this trick we will eventually get to
above 1\% but it's definitely not the most practical solution.}
\vspace{0.2cm}

\textit{It's natural to think about how far I can push this trick, and
whether or not I can do this trick generally. Actually since our initial
discovery gave us roughly halfway progress, I wondered if I could
"flip" the signs of that equation to get another $+ 0.5\%$ progress. More
precisely, the linear diophantine equation $27x + 8y$ also admits the
solution}
$$
27(-5) + 8(17) = 1.
$$
\textit{From this point we can use the same tricks as before to get that
numbers of the form $27(5 + 8n)$ are mischievous. I was right this class
of numbers gives another $1/216 \approx 0.46 \%$ to the density of
mischievous numbers. Finally to get that final $\approx 0.08\%$ I can
just use the numbers $125(3+8n)$. This should be the problem solved,
I'll write everything up cleanly tomorrow.}
\end{hint}
\newpage

\begin{solution}{Solution}
We will show that there are certain arithmetic progressions such that
every term in those progressions is mischievous. Then we will show that
these numbers appear often enough to be able to guarentee that at least
$\approx 1.02 \%$ of natural numbers are mischievous. \vspace{0.2cm}

We begin with a couple of claims.
\vspace{0.2cm}

\textit{Claim 1: Every number of the form $8(-17 + 27n)$ is mischievous
for $n \in \mathbb{N}$.}
\vspace{0.2cm}

\textit{Proof:} Let $k_n = 8(-17 + 27n)$, since $8 \mid k_n$, $k_n$ is
divisible by a perfect cube. $k_n + 1 = 216n - 135 = 27(8n - 5)$ and so
$27 \mid (k_n + 1)$ and $(k_n + 1)$ is divisible by a perfect cube too.
Thus $k_n$ is always mischievous. $\square$
\\

\textit{Claim 2: Every number of the form $27(-3 + 8n)$ is mischievous
for $n \in \mathbb{N}$.}
\vspace{0.2cm}

\textit{Proof:} Let $k_n = 27(-3 + 8n)$, since $27 \mid k_n$, $k_n$ is
divisible by a perfect cube. $k_n + 1 = 216n - 80 = 8(27n - 10)$ and so
$8 \mid (k_n + 1)$ and $(k_n + 1)$ is divisible by a perfect cube too.
Thus $k_n$ is always mischievous. $\square$
\\

\textit{Claim 3: Every number of the form $125(-5 + 8n)$ is mischievous
for $n \in \mathbb{N}$.}
\vspace{0.2cm}

\textit{Proof:} Let $k_n = 125(-5 + 8n)$, since $125 \mid k_n$,
$k_n$ is divisible by a perfect cube. $k_n + 1 = 1000n - 624 =
8(125n - 78)$ and so $8 \mid (k_n + 1)$ and $(k_n + 1)$ is divisible
by a perfect cube too. Hence $k_n$ is always mischievous. $\square$ \\

\bigskip
Define three sequences given by
$$
a_n = 8(-17 + 27n), \quad b_n = 27(-3 + 8n), \quad c_n = 125(-5 + 8n)
$$
\bigskip

\textit{Claim 4: For any set $A$ of consecutive natural numbers of size
$N \geq 27(8) = 216$, there must be at least one term from $a_n$ and
at least one term from $b_n$ in the set $A$.} \vspace{0.2cm}

\textit{Proof:} For both sequences it is easy to verify that
$$
a_{n+1} - a_n = 216, \quad b_{n+1} - b_n = 216.
$$
I will prove the following auxilliary fact. Let $x \in \mathbb{N}$.
If $x \equiv 80 \ (\text{mod } 216)$ then $x$ is a part of $(a_n)$.
Note that $x \equiv 80 - 216 \equiv -136 \ (\text{mod } 216)$. So we may
write $x = 216n - 136 = 8(-17 + 27n)$, where $n \in \mathbb{N}$ because
if $n < 1$ then this would imply that $x < 0$ which is impossible. Thus
$x  = a_n$ is a part of the sequence $(a_n)$. \vspace{0.2cm}

Now among a set of $k \geq 216$ consecutive numbers there must be all
residues modulo 216. Thus there is a number $\equiv 80$ modulo 216
and this number is a part of $(a_n)$. \vspace{0.2cm}

We again prove an auxilliary result. Let $x \in \mathbb{N}$. If
$x \equiv 135 \ (\text{mod } 216)$, then $x$ is a part of $(b_n)$. Observe
that $x \equiv 135 - 216 \equiv -81 \ (\text{mod } 216)$. Thus we can write
$x = 216n - 81 = 27(-3 + 8n)$ where $n \in \mathbb{N}$ because if $n < 1$
then this would imply that $x < 0$ which isn't possible. Hence
$x = b_n$ and is a part of the sequence $(b_n)$. \vspace{-0.2cm}

Now again among a set of $k \geq 216$ consecutive numbers there must be
all residues modulo 216. Thus there must be a number $\equiv 135$
modulo 216 and this number is a part of $(b_n)$.
$\square$ \\ \bigskip

\textit{Claim 5: For any set $A$ of consecutive natural numbers of size
$N \geq 125(8) = 1000$, there must be at least one term from $c_n$ in the
set $A$.}
\vspace{0.2cm}

\textit{Proof:} Indeed one may check that
$$
c_{n+1} - c_n = 1000.
$$
Let's show the following auxilliary result. Let $x \in \mathbb{N}$. If
$x \equiv 375 \ (\text{mod } 1000)$, then $x$ is a part of $(c_n)$. Observe
that $x \equiv 375 - 1000 \equiv -625 \ (\text{mod } 1000)$. So we can write
$x = 1000n - 625 = 125(-5 + 8n)$ for some $n \in \mathbb{N}$ because
otherwise if $n < 1$ then $x < 0$ which is impossible. Thus $x = c_n$ and
is part of the sequence $(c_n)$.
\vspace{0.2cm}

Now among a set of $k \geq 1000$ consecutive numbers there must be all
residues modulo 1000. Thus there is a number that is $\equiv 375$ mod 1000,
and this number is a part of $(c_n)$. $\square$ \\
\bigskip

\textit{Claim 6: The following is true,}
\[
\begin{aligned}
\{a_1, a_2, \ldots \} \cap \{b_1, b_2, \ldots \} &= \emptyset\\
\{a_1, a_2, \ldots \} \cap \{c_1, c_2, \ldots \} &= \emptyset\\
\end{aligned}
\]

\textit{Proof:} A number cannot be part of both $(a_n)$ and $(b_n)$,
since it would
imply that the number is congruent to both 135, 80 modulo 216 which is
impossible. Suppose some number $k$ is a part of both $(a_n)$ and $(c_n)$.
Then $k = 1000m - 625$ for some $m \in \mathbb{N}$ and we require
\[
\begin{aligned}
1000m - 625 &\equiv 80 \ (\text{mod } 216) \\
136n &\equiv 57 \ (\text{mod } 216).
\end{aligned}
\]
Now if there is a solution to this congruence then there must exist
$x, y \in \mathbb{Z}$ such that $136x + 216y = 57$. Bézout's theorem tells
us that a solution will only exist if $\gcd(136, 216) = 8$ divides 57 which
isn't true. $\square$
\\ \bigskip

\textit{Claim 7: Let $k = 1000m - 625 = c_m, (m\in \mathbb{N})$ be a part of
$(c_n)$. Then $k$ is a part of $(b_n)$ if and only if
$$
m \in \{ 10,37,64,91,118,145,172,199 \} \bmod 216.
$$
}

\textit{Proof:} For $k$ to be a part of $(b_n)$ we need,
\[
\begin{aligned}
1000m - 625 &\equiv 135 \ (\text{mod } 216) \\
136m &\equiv 112 \ (\text{mod } 216) \\
\end{aligned}
\]
It suffices to study the linear diophantine equation
\[
\begin{aligned}
136m + 216y &= 112\\
17m + 27y &= 14 \\
\end{aligned}
\]
Now using whatever method you like, say Euclidean algorithm backtracking
you can find that $17^{-1} \equiv 20 \ (\text{mod } 27)$ and thus
$m \equiv 14 \cdot 20 \equiv 10 \ (\text{mod } 27)$. Now
in mod 216 the only such $m$ are
$$
m \in \{ 10,37,64,91,118,145,172,199 \}. \ \square
$$
\\

\textit{Claim 8: In any group of 27 consecutive terms from $(c_n)$
there is exactly one term from $(b_n)$.} \vspace{0.2cm}

\textit{Proof:} Terms from $(c_n)$ take on the form $1000n - 625$.
With 27 consecutive terms, $n$ takes on all residues modulo 27 exactly once.
Now due to claim 7, a number is a part of $(b_n)$ if $n \equiv 10
\ (\text{mod } 27)$. Among these 27 consecutive terms exactly one
has $n \equiv 10 \ (\text{mod } 27)$ so exactly one term is a part of
$(b_n)$. $\square$
\\ \bigskip

Now let's take any set of $125(8)(27) = 27000$ consecutive natural numbers.
\begin{itemize}
\item In this set there are $\frac{27000}{216} = 125$ terms from $(a_n)$ and
$\frac{27000}{216} = 125$ terms from $(b_n)$, due to claim 4.
\item In this set there are $\frac{27000}{1000} = 27$ terms from $(c_n)$,
due to claim 5.
\end{itemize}
The numbers from these three sequences are all mischievous by claims 1,2,3.
We know that the only times that terms can be a part of multiple sequences
is when a term is a part of both $(b_n)$ and $(c_n)$ due to claim 6.
Due to claim 7, we know that among those 27 terms from $(c_n)$, there is
exactly one that is also a part of $(b_n)$, so we must take this into
account. \vspace{0.2cm}

Thus in any set of 27000 consecutive natural numbers there are at least
$125 + 125 + 27 - 1 = 276$ mischievous numbers.
Hence among these 27000 consecutive natural
numbers, at least
$$
\frac{276}{27000} = 1.0 \dot{2} \% > 1\%
$$
of the numbers are mischievous. $\blacksquare$
\\

\textit{Comments: Really enjoyed working through this one!}
\\

\hline
\vspace{0.2cm}

\textit{Addendums (after having read solutions):} Problem 7 was a really
nice one. I guess this is an analysis problem? I essentially found the
same solution as the official here, actually I think our constructions are
almost exactly the same, the only thing is I didn't include numbers
$\equiv 624 \ (\text{mod } 1000)$, I think that's why my final density
ended up dangerously close to $1 \%$.

\end{solution}
\newpage

\begin{problem}{Problem 8}
(a) Historically, it used to be possible to buy chicken McNuggets in boxes
of 9 or 20. What was the largest number of chicken McNuggets that it was
impossible to buy. \vspace{0.2cm}

(b) If instead they were in boxes of size $r$ and size $s$, with
$r$ and $s$ coprime, show that the largest number that it is impossible
to buy is $rs - r - s$. \textit{This is actually called the Chicken
McNugget Theorem!}
\vspace{0.2cm}

\textit{You might want to consider which numbers are possible for each
value modulo $s$, if $s > r$.}
\vspace{0.2cm}

(c) \textit{(Extension)} Show that there are exactly
$\frac12 (r-1)(s-1)$ positive integers that it is impossible to purchase.
\end{problem}

\begin{solution}{Solution}
(a) Say that we can express a natural number $n$ \textit{nicely} if there
exists a solution to $n = 9x + 20y$, for $x,y \in \mathbb{Z}_{\geq 0}$.

Let $x,y$ be non-negative integers. It is possible to buy
$9x + 20y$ chicken McNuggets. Bézout's theorem tells us that for
$x, y \in \mathbb{Z}$ there is always a solution to
$n = 9x + 20y$ since $\gcd(9, 20) = 1$. So for it to be "impossible" to
buy some amount of chicken McNuggets, there must be some condition
which forces one of $x, y < 0$. \vspace{0.2cm}

\textit{Claim 1: Let $n = 9x + 20y$ be positive with $x,y \in \mathbb{Z}$.
If $y \geq 9$ then we can always express $n$ nicely.} \vspace{0.2cm}

\textit{Proof:} If $x \geq 0$ then there is nothing to prove so
assume that $x < 0$. For convenience let us write $a = -x > 0$.
We can induct on $y$. \vspace{0.2cm}

For the base case, note that if $n = 20(9) - 9a$, then $a < 20$ since
$n > 0$. Now $n = -9(a - 20) = 9(20-a)$, and since $20-a > 0$ the base
case is true. \vspace{0.2cm}

Now assume the result is true for $y = k \geq 9$. That means that if
$n = 20k - 9a$, then we can write it as $n = 9x' + 20y'$ for some
$x', y' \in \mathbb{Z}_{\geq 0}$. Now consider,
\[
\begin{aligned}
m &= 20(k+1) - 9a \\
&= (20k - 9a) + 20 \\
&= (9x' + 20y') + 20 \\
&= 9x' + 20(y' + 1)
\end{aligned}
\]
and so the result is also true for $k+1$. Thus the proof is complete
via induction. $\square$
\\

Suppose that $N = 9x + 20y$ is the largest number of nuggets that
it's impossible to buy. Due to claim 1, we know that $y \leq 8$. Since we
want to find the largest such $N$ we should just pick $y = 8$. Now we
have $N = 20(8) + 9x$. Clearly if $x \geq 0$ then it is possible to buy
$N$ nuggets so we must have $x < 0$. Again since we want to maximise $N$
it is best to pick $x = -1$, thus $N = 20(8) - 9(1) = 151$. \vspace{0.2cm}

To finish we need to show that it isn't possible to express $N$ nicely.
Suppose $151 = 9x + 20y$ is written nicely. Taking mod 9 this becomes
$2y \equiv 7 \ (\text{mod } 9)$ from which you can use, say Bézout's
theorem to show that $y \equiv 8 \ (\text{mod } 9)$. If $y \geq 9$ then
this forces $x < 0$, so this isn't possible. Also $y \geq 0$ by assumption
so the only candidate for $y$ is $y = 8$, from which we can solve to get
$x = -1$, a contradiction. $\square$ \\ \bigskip

(b) The general case is similar. We have $s,r \in \mathbb{N}$ with
$\gcd(r,s) = 1$. Let us assume $s > r$. \vspace{0.2cm}

\textit{Claim 1: Let $n = sX + rY$ be positive with $X,Y \in \mathbb{Z}$.
If $X \geq r$ then we can express $n$ nicely.} \vspace{0.2cm}

\textit{Proof: } If $Y \geq 0$ then there is nothing to prove. So let us
assume $Y < 0$. For convenience we will write $A = -Y > 0$.
We are going to induct on $X$. \vspace{0.2cm}

For the base case, we see that if $n = s(r) - rA$, then $A < s$ since
$n$ is positive. Now $n = r(s - A)$ and since $(s - A) > 0$ the base
case is indeed true. \vspace{0.2cm}

Now assume the result holds for $X = k \geq r$. Thus if $n = sk - rA$ we
can write it as $n = sX' + rY'$ for some $X', Y' \in \mathbb{Z}_{\geq 0}$.
Consider,
\[
\begin{aligned}
m &= s(k + 1) - rA \\
&= (sk - rA) + s \\
&= (sX' + rY') + s \\
&= s(X' + 1) + rY'
\end{aligned}
\]
and so the result is true for $X = k+1$ also. Thus the proof is complete
via induction. $\square$ \\

Suppose that $N = sX + rY$ is the largest number of nuggets that it's
impossible to buy. Now due to claim 1, we must have $X < r$ and since we
want to maximise $N$ we will just pick $X = r-1$. We have
$N = s(r-1) + rY$. Obviously if $Y \geq 0$ then $N$ has been written nicely
so $Y < 0$. Since we want to maximise $N$ we will pick $Y = -1$.
So we have $N = s(r-1) - r = rs - s -r$. \vspace{0.2cm}

To finish we need to show that it isn't possible to express $N$ nicely.
Suppose FTSOC that
$$
N = rs - s - r = sX + rY
$$
has been written nicely. Taking mod $r$ we have
\[
\begin{aligned}
-s &\equiv sX \ (\text{mod } r) \\
-1 &\equiv X \ (\text{mod } r) && (*) \\
\end{aligned}
\]
\textit{($*$ --- $s$ has a modular inverse mod $r$ because $\gcd(r,s) = 1$)}
. If
$X \geq r$ then the RHS will become too large and this will force $Y < 0$,
so $X < r$. Also $X \geq 0$ since $N$ was written nicely. Since
$X \equiv -1 \ (\text{mod } r)$, these two inequalities mean that the only
candidate for $X$ is $X = r-1$. But now plugging this in and solving
yields $Y = -1$, a contradiction. $\blacksquare$ \\

\textit{Remark: The solution
I found is much nicer to follow if you consider mod $r$ instead of mod $s$.}
\vspace{0.2cm}

\textit{
I'll write a little more detail about $(*)$. Let $x$ be the modular
inverse of $a$ mod $b$. Then $ax \equiv 1 \ (\text{mod } b)$. This means
that we can find $x,y \in \mathbb{Z}$ such that $ax + by = 1$. From
Bézout's theorem we know that solutions only exist if $\gcd(a,b) \mid 1$,
that is $\gcd(a,b) = 1$.
}

\newpage
(c) We'll take the same setup as in part (b). That is let
$s,r \in \mathbb{N}$ with $\gcd(r,s) = 1$ and $s > r$.
Similar to claim
1, in part (b) it is natural to guess that the following result is also
true. \\

\textit{Claim 1: Let $n = sX + rY$ be positive with $X, Y \in \mathbb{Z}$.
If $Y \geq s$ then we can express $n$ nicely.}
\vspace{0.2cm}

\textit{Proof:} If $X \geq 0$ there is nothing to prove, so assume
$X < 0$. For convenience let $A = -X > 0$. We will induct on $Y$.
\vspace{0.2cm}

For the base case when $Y = s$, we have $n = r(s) - sA$ with $A < r$ since
$n > 0$. Now $n = s(r - A)$ and since $r-A > 0$ the base case is true.
\vspace{0.2cm}

Now assume that the result holds for $Y = k \geq s$. So if
$n = rk - sA$ we can write it as $n = sX' + rY'$ for some
$X', Y' \in \mathbb{Z}_{\geq 0}$. Consider
\[
\begin{aligned}
m &= r(k+1) - sA \\
&= (rk - sA) + r \\
&= (sX' + rY') + r \\
&= sX' + r(Y' + 1)
\end{aligned}
\]
and so the result is also true for $k+1$. The proof is finished via
induction. $\square$ \\

Let $n = sX + rY$ be positive and impossible to write nicely.
Exactly one of $X, Y$ must be negative. This is because if they were both
negative than this would mean that $n < 0$, and if they were both $\geq 0$
then $n$ has been written nicely. \\

\textit{Claim 2: For the $n$ described above, we can assume $Y < 0$.}
\vspace{0.2cm}

\textit{Proof:} We will show that any such $n$ with $X < 0$ can be
re-written in such a way that $Y < 0$. Suppose indeed that $Y > 0, X < 0$,
let $A = -X > 0$ for convenience. So we have $n = rY - sA$ being positive
and impossible to write nicely. Due to claim 1 (c) we must have $Y < s$.
Also since
$$
0 < n = rY - sA < rs - sA
$$
we have $sA < rs \Rightarrow A < r$.
\[
\begin{aligned}
n &= rY - sA \\
&= rY - sA + (rs - rs) \\
&= r(Y - s) - s(A - r) \\
&= s(r - A) + r(Y - s)
\end{aligned}
\]
but thanks to the work above we know that $Y - s < 0$ and $r - A > 0$,
so if we just let $X' = r-A$ and $Y' = Y-s$, we have written
$n = sX' + rY'$ with $X' > 0, Y' < 0$. $\square$ \\

Now thanks to claim 2, we safely assume $X > 0$ and $Y < 0$. Due to
claim 1 (b) we have
$$
X \in \{ 1,2,\ldots,r-1 \}
$$
Let $B = -Y > 0$ for convenience. We want to find the set of values that
$B$ can take on. Since $n > 0$,
\[
\begin{aligned}
n = sX - rB &> 0 \\
sX &> rB \\
\frac{sX}{r} &> B \\
B &\leq \floor{ \frac{sX}{r} }
\end{aligned}
\]
where the last line is because $B \in \mathbb{Z}$. So for each selection
of $X$ there are $\floor{sX/r}$ choices for $B$ or equivalently
$$
B \in \left\{ 1,2,\ldots, \floor{\frac{sX}{r}} \right\}.
$$
Thus putting this together there are precisely
$$
S = \sum_{X=1}^{r-1} \floor{ \frac{sX}{r} }
$$
integers that are impossible to write nicely. So now our goal is to
evaluate $S$. Though it looks like a nasty sum, it is possible to
evaluate it nicely. We need a key lemma. \\

\textit{Claim 3: For $X = 1,2,\ldots, r-1$, the number $sX$ leaves
all non-zero residues modulo $r$. Precisely that is
$$
\{ s(1), s(2),\ldots, s(r-1) \} = \{ 1, 2, \ldots, r-1 \}.
$$
where all those numbers are reduced modulo $r$.
}
\vspace{0.2cm}

\textit{Proof:} Suppose FTSOC that there is some integer $1\leq q \leq r-1$,
such that there does not exist $X = 1, 2,\ldots,r-1$ with
$$
sX \equiv q \ (\text{mod } r)
$$
Since $\gcd(s,r) = 1$ we can write
$$
X \equiv s^{-1} q \ (\text{mod } r)
$$
Now the only way that $X$ cannot exist in this scenario is if
$s^{-1} q \equiv 0$. However $q > 0$ and the modular inverse of $s$ is
obviously not 0, so we have reached a contradiction. $\square$ \\

Now we can write
\[
\begin{aligned}
S &= \sum_{X=1}^{r-1} \frac{sX}{r} - \left\{ \frac{sX}{r} \right\} \\
&= \frac{s}{r} \frac{(r-1)r}{2} -
\sum_{X=1}^{r-1} \left\{ \frac{sX}{r} \right\} \\
&= \frac{s(r-1)}{2} -
\sum_{X=1}^{r-1} \left\{ \frac{sX}{r} \right\} \\
\end{aligned}
\]
Now thanks to claim 3, $sX$ leaves every non-zero residue mod $r$.
Since the fractional part of $\frac{sX}{r}$ is just
$\frac{sX \bmod r}{r}$ we can write
\[
\begin{aligned}
&= \frac{s(r-1)}{2} -
\sum_{k=1}^{r-1} \frac{k}{r} \\
&= \frac{s(r-1)}{2} - \frac{1}{r} \frac{(r-1)r}{2} \\
&= \frac{s(r-1)}{2} - \frac{(r-1)}{2} \\
&= \frac12 (s-1)(r-1). \ \blacksquare
\end{aligned}
\]

\textit{Comments: This one was really nice also, lots of moving parts here.
Throughout that solution $\floor{\cdot}$ refers to the floor of a number,
and $\{ \cdot \}$ refers to the fractional part of a number, (or the
brackets enclose a set, it is dependent on context.)}
\\

\hline
\vspace{0.2cm}

\textit{Addendums (after having read solutions):} My solutions (once again)
look very different to the official ones. It seems that the official
solutions are much more succint and avoid using much technical language,
which is understandable since I guess UKMT wants the solutions to be
understandable even to those who have only the tools at GCSE maths level.
I think a nice way to describe it is that the UKMT solutions
are more "philosophical" as in they reason with pure logic why the answer
should be something, whilst my solutions are more explicit or computational
showing why an answer should be true with concrete calculations and
theorems.
\vspace{0.2cm}

Though my solution is much longer, I particularly liked my solution to
part (c) here. At first it seems really impossible to evaluate $S$
as it doesn't look like any of the standard summations that we have learn't
to deal with, but a nice solution does in fact exist using the floor and
fractional part functions.
\end{solution}

\end{document}
